% concatenated from 201214647.tex


\documentclass{article}
\usepackage{amsmath, amssymb, amsthm, graphicx, bm, url,color}
  %,comment}
  % cite,
%\usepackage[all]{xy}


%\usepackage{tikz}
%\usetikzlibrary{matrix,arrows,calc,intersections,fit}
%\usepackage{tikz-cd}

%%%
%%%%%%%%%%%%%%%%%%%%%%%%%%%%%%%%%%%%%%%%%%%%%%%%%%%%%%%%%%%%%%%%%%%%%%%%
\DeclareMathOperator{\Mod}{mod}
\DeclareMathOperator{\supp}{supp}
\DeclareMathOperator{\res}{res}
\DeclareMathOperator{\grad}{grad} \DeclareMathOperator{\Hom}{Hom}
\DeclareMathOperator{\cone}{Cone} \DeclareMathOperator{\hess}{hess}
\DeclareMathOperator{\gr}{graph} \DeclareMathOperator{\tr}{tr}
\DeclareMathOperator{\phase}{phase} \DeclareMathOperator{\di}{div}
\DeclareMathOperator{\area}{area} \DeclareMathOperator{\re}{Re}
\DeclareMathOperator{\im}{im} \DeclareMathOperator{\spt}{Spt}
\DeclareMathOperator{\coker}{coker}
\DeclareMathOperator{\crit}{Crit} \DeclareMathOperator{\obj}{obj}
\DeclareMathOperator{\Ind}{Ind}
\DeclareMathOperator{\Ham}{Ham}
\DeclareMathOperator{\Lag}{Lag}
\DeclareMathOperator{\Exp}{Exp}
\DeclareMathOperator{\sing}{Sing}
\DeclareMathOperator{\ev}{ev}

\DeclareMathOperator{\rel}{rel}
\DeclareMathOperator{\val}{val}
\DeclareMathOperator{\Spin}{Spin}
\DeclareMathOperator{\Loc}{Loc}
\DeclareMathOperator{\Def}{Def}
\DeclareMathOperator{\LG}{LG}

%%%%%%%%%%%%%%%%%%%%%%%%%%%%%%%%%%%%%%%%%%%%%%%%%%%%%%%%%%%%%%%%%%%%%%%%
%%%%%%%%%%%%%%%%%%%%%%%%%%     Macros      %%%%%%%%%%%%%%%%%%%%%%%%%%%%%
%%%%%%%%%%%%%%%%%%%%%%%%%%%%%%%%%%%%%%%%%%%%%%%%%%%%%%%%%%%%%%%%%%%%%%%%
\theoremstyle{plain}% default
\newtheorem{thm}{Theorem}[section]
\newtheorem{prop}[thm]{Proposition}
\newtheorem{lem}[thm]{Lemma}
\newtheorem{sblem}[thm]{SubLemma}
\newtheorem{cor}[thm]{Corollary}
\theoremstyle{definition}
\newtheorem{ax}[thm]{Axiom}
\newtheorem{nt}[thm]{Notation}\newtheorem{conv}[thm]{Convention}
\newtheorem{cj}[thm]{Conjecture}
\newtheorem{qs}[thm]{Question}
\newtheorem{dfn}[thm]{Definition}
\newtheorem{prb}[thm]{Problem}
\newtheorem{prn}[thm]{Principle}
\newtheorem{pr}[thm]{Property}
\newtheorem{cn}[thm]{Construction}
\newtheorem{cond}[thm]{Condition}
\newtheorem{exam}[thm]{Example}
\newtheorem{rem}[thm]{Remark}
\newtheorem{hp}[thm]{Hypothesis}
\newtheorem{stu}[thm]{Situation}
\newtheorem*{ackn}{Acknowledgements}


\numberwithin{equation}{section}
\numberwithin{figure}{section}


\def\e#1\e{\begin{equation}#1\end{equation}}
\def\iz#1\iz{\begin{itemize}#1\end{itemize}}
\def\ea#1\ea{\begin{align}#1\end{align}}
\def\eq{\eqref}
\def\l{\label}
\def\0{\hspace{0pt}}
\def\cyl{\rm{cyl}}
\def\tW{\,\widetilde{\!\cW}}
\def\alti{\widetilde\al}
\def\Bti{\,\widetilde{\!B}}
\def\Eti{\,\widetilde{\!E}}
\def\Pti{\,\widetilde{\!P}}
\def\Qti{\,\widetilde{\!Q}}
\def\Qh{\,\widehat{\!Q}}
\def\Tti{\,\widetilde{\!T}}
\def\Thti{\widetilde{\Th}}
\def\thti{\ti\th}
\def\Phiti{\widetilde{\Phi}}
\def\Lati{\widetilde{\La}}
\def\phiti{\widetilde{\phi}}
\def\xiti{\widetilde{\xi}}
\def\Upti{\widetilde{\Up}}
\def\Dti{\,\widetilde{\!D}}
\def\Lti{\,\widetilde{\!L}}
\def\Lh{\,\widehat{\!L}}
\def\Nti{\widetilde N}
\def\Sti{\widetilde S}
\def\Uti{\widetilde{U}}
\def\Uh_#1{\,\widehat{\!U}_{\!#1}}
\def\Vti{\,\widetilde{\!V}}
\def\VtiKx{\,\widetilde{\!V}_{\!K,x}}
\def\Vtix{\,\widetilde{\!V}_{\!x}}
\def\Fti{\widetilde\F}
\def\Gti{\widetilde{\rm G}}
\def\Wh{\widehat{W\,}\!_{\!x}}
\def\Wtix{\widetilde{W\,}\!_{\!x}}
\def\Wti{\widetilde{W\,}\!}
\def\Xti{\widetilde X}
\def\Xh{\widehat X}
\def\Bh{\widehat B}
\def\Kh{\widehat K}
\def\Yh{\widehat Y}
\def\Yti{\widetilde Y}
\def\Ztix{\,\widetilde{\!Z}_{\!x}}
\def\Zhx{\,\widehat{\!Z}_{\!x}}
\def\Zh{\,\widehat{\!Z}}
\def\Jh{\hat J}
\def\Xih{\widehat\Xi}
\def\Omh{\widehat\Om}
\def\gh{\hat g}
\def\gti{\ti g}
\def\hh{\hat h}
\def\omh{\widehat\om}
\def\omti{\widetilde\om}
\def\Omti{\widetilde\Om}
\def\Gati{\widetilde\Ga}
\def\gati{\widetilde\ga}
\def\Psiti{\widetilde\Psi}
\def\rhoti{\widetilde\rho}
\def\Xiti{\widetilde\Xi}
\def\psiti{\widetilde{\psi}}
\def\etati{\widetilde{\eta}}
\def\Io{{\rm I}}
\def\fa{\mathfrak a}
\def\fb{\mathfrak b}
\def\fg{\mathfrak g}
\def\fm{\mathfrak m}
\def\fp{\mathfrak p}
\def\fsu{\mathfrak{su}}
\def\px{\approx}
\def\F{{\rm F}}
\def\str{{\rm str}}
\def\nv{{\rm nv}}
\def\nd{{\rm nd}}
\def\rs{{\rm rs}}

\def\T*#1{T^{*\!\!}#1}

\def\Ind{\mathop{\rm Ind}\nolimits}
\def\Fix{\mathop{\rm Fix}\nolimits}
\def\Curv{\mathop{\rm Curv}\nolimits}
\def\Res{\mathop{\rm Res}\nolimits}
\def\Bl{\mathop{\rm Bl}\nolimits}
\def\dim{\mathop{\rm dim}\nolimits}
\def\det{\mathop{\rm det}\nolimits}
\def\rank{\mathop{\rm rank}}
\def\id{\mathop{\rm id}\nolimits}
\def\GL{\mathop{\rm GL}}
\def\SO{\mathop{\rm SO}\nolimits}
\def\Ker{\mathop{\rm Ker}}
\def\CoKer{\mathop{\rm CoKer}}
\def\Re{\mathop{\rm Re}}
\def\Im{\mathop{\rm Im}}
\def\Aut{\mathop{\rm Aut}}
\def\End{\mathop{\rm End}}
\def\SU{\mathop{\rm SU}}
\def\U{{\rm U}}
\def\SLag{\mathop{\rm SLag}}
\def\Hess{\mathop{\rm Hess}}
\def\Hom{\mathop{\rm Hom}\nolimits}
\def\vol{\mathop{\rm vol}\nolimits}
\def\an{{\rm an}}
\def\ge{\geqslant}
\def\le{\leqslant\nobreak}
\def\les{\lesssim\nobreak}
\def\ges{\gtrsim\nobreak}

\def\B_#1{\,\ov{\!B_{#1}}}
\def\Ba_#1{\,\ov{\!B_{#1}\!}\,}

\def\cA{{\mathbin{\cal A}}}
\def\cB{{\mathbin{\cal B}}}
\def\cC{{\mathbin{\cal C}}}
\def\cD{{\mathbin{\cal D}}}
\def\cE{{\mathbin{\cal E}}}
\def\cF{{\mathbin{\cal F}}}
\def\cG{{\mathbin{\cal G}}}
\def\cH{{\mathbin{\cal H}}}
\def\cI{{\mathbin{\cal I}}}
\def\cJ{{\mathbin{\cal J}}}
\def\cK{{\mathbin{\cal K}}}
\def\cL{{\mathbin{\cal L}}}
\def\cM{{\mathbin{\cal M}}}
\def\Mbar{\overline{\cM}}
\def\bM{\,\ov{\!\cM}}
\def\bN{\,\ov{\!\cN}}
\def\cN{{\mathbin{\cal N}}}
\def\oM{{\mathbin{\smash{\,\,\overline{\!\!\mathcal
M\!}\,}}}}
\def\cO{{\mathbin{\cal O}}}
\def\cP{{\mathbin{\cal P}}}
\def\cR{{\mathbin{\cal R}}}
\def\cS{{\mathbin{\cal S}}}
\def\cT{{\mathbin{\cal T}}}
\def\cU{{\mathbin{\cal U}}}
\def\cV{{\mathbin{\cal V}}}
\def\={\equiv}

\def\PV{{\mathbin{\cal{PV}}}}
\def\TS{{\mathbin{\cal{TS}}}}
\def\cQ{{\mathbin{\cal Q}}}
\def\cW{{\mathbin{\cal W}}}
\def\cX{{\mathbin{\cal X}}}
\def\cY{{\mathbin{\cal Y}}}
\def\cZ{{\mathbin{\cal Z}}}
\def\CY{{\mathbin{\cal CY}}}
\def\SL{{\mathbin{\cal SL}}}


\def\C{{\mathbin{\mathbb C}}}
\def\CP{{\mathbin{\mathbb{C}}}P}
\def\RP{{\mathbin{\mathbb{R}}}P}
\def\OO{\O(-1)^{\op2}}

\def\K{{\mathbin{\mathbb K}}}
\def\Q{{\mathbin{\mathbb Q}}}
\def\R{{\mathbin{\mathbb R}}}
\def\Z{{\mathbin{\mathbb Z}}}
\def\al{\alpha}
\def\be{\beta}
\def\ga{\gamma}
\def\de{\delta}
\def\et{\eta}
\def\io{\iota}
\def\ep{\epsilon}
\def\la{\lambda}
\def\ka{\kappa}
\def\th{\theta}
\def\ze{\zeta}
\def\up{\upsilon}
\def\vp{\varphi}
\def\si{\sigma}
\def\om{\omega}
\def\ph{\phi}
\def\Ph{\Phi}
\def\ps{\psi}
\def\Ps{\Psi}
\def\r{\rho}
\def\ta{\tau}
\def\De{\Delta}
\def\La{\Lambda}
\def\Si{\Sigma}
\def\Tau{{\rm T}}
\def\Th{\Theta}
\def\Om{\Omega}
\def\Ga{\Gamma}
\def\Up{\Upsilon}
\def\sSi{{\smash{\sst\Si}}}

\def\bmu{\bm\mu}
\def\bnu{\bm\nu}
\def\bw{\bigwedge}
\def\Id{{\rm Id}}
\def\fl{{\rm flat}}
\def\nv{{\rm nov}}

\def\d{{\partial}}
\def\db{\bar\d}%{\,\ov{\!\partial}}
\def\ts{\textstyle}
\def\st{\scriptstyle}
\def\sst{\scriptscriptstyle}
\def\w{\wedge}
\def\-{\setminus}
\def\bl{\bullet}
\def\op{\oplus}
\def\bop{\bigoplus}
\def\ot{\otimes}
\def\bot{\bigotimes}
\def\ov{\overline}
\def\ovr{\overrightarrow}
\def\ul{\underline}
\def\iy{\infty}
\def\es{\emptyset}
\def\ra{\rightarrow}
\def\Ra{\Rightarrow}
\def\Longra{\Longrightarrow}
\def\ab{\allowbreak}
\def\longra{\longrightarrow}
\def\hookra{\hookrightarrow}
\def\dashra{\dashrightarrow}
\def\t{\times}
\def\cm{\circ}
\def\ti{\tilde}
\def\hf{{\frac{1}{2}}}
%\def\md#1{\vert #1 \vert}
\def\bmd#1{\big\vert #1 \big\vert}
\def\ms#1{\vert #1 \vert^2}
\def\nbk{\nabla_{\!k\,}}
\def\nbl{\nabla_{\!l\,}}
\def\nbi{\nabla_{\!i\,}}
\def\nb{\nabla}
\def\sb{\subseteq}
\def\k{{\mathbf{k}}}
\def\b{{\mathbf{b}}}
\def\x{{\mathbf{x}}}
\def\bfb{{\mathbf{b}}}
\def\v{{\mathbf{v}}}
\def\ig{{   {\rm int} }}
\def\ex{{   {\rm ex}    }}
\def\curv{{   {\rm curv}    }}
\def\m{{   {\mathfrak m}    }}
%%%%%%%%%%%%%%%%%%%%%% Text of paper %%%%%%%%%%%%%%%%%%%%%%%%%%%%%%%%%%%
%%%%%%%%%%%%%%%%%%%%%%%%%%%%%%%%%%%%%%%%%%%%%%%%%%%%%%%%%%%%%%%%%%%%%%%%
\title{Generalized Thomas--Yau Uniqueness Theorems}
\author{Yohsuke Imagi}
\begin{document}
\maketitle

%\tableofcontents

\begin{abstract}
We generalize Thomas--Yau's uniqueness theorem \cite[Theorem 4.3]{TY} in two ways.
We prove a stronger statement for special Lagranigans and
include minimal Lagrangians in K\"ahler--Einstein manifolds
or more generally $J$-minimal Lagrangians introduced by Lotay and Pacini \cite{LP,LP2}.
In every case the heart of the proof is to make certain Hamiltonian perturbations.
For this we use the method by Imagi, Joyce and Oliveira dos Santos \cite[Theorem 4.7]{IJO}.
%We also generalize Fukaya, Oh, Ohta and Ono's $HF^*$ nonvanishing theorem \cite[Theorem E]{FOOO}
%to get the uniqueness results.
%For this we do not follow Thomas--Yau's original method but use that in
%the author's previous joint paper with Joyce and Oliveira dos Santos \cite{IJO}.
%It relies upon \L{}ojasiewicz's estimate for real analytic functions,
%but by a suitable perturbation argument we can make the results hold
%in the $C^\iy$ category.
\end{abstract}
\section{Introduction}
In this paper we improve and generalize Thomas--Yau's theorem \cite[Theorem 4.3]{TY}.
Our first main result is the following.
For the more complete statement see Corollary \ref{2} (ii).
\begin{thm}\l{01}
Let $X$ be a K\"ahler manifold equipped with a holomorphic volume form.
Let the cohomology Fukaya category $H\cF(X)$ have two isomorphic objects supported near two closed irreducibly-immersed special Lagrangians $L_1,L_2$ respectively.
Then $L_1=L_2\sb X.$
\end{thm}
\begin{rem}
{\bf(i)}
Thomas and Yau prove their uniqueness theorem for closed special Lagrangians $L_1,L_2$ in the same Hamiltonian isotopy class.
But their proof works under the weaker hypothesis as above; that is, we need only that $L_1,L_2$ have the same isomorphism class in $H\cF(X).$
This $H\cF(X)$ is a (usual) associative category obtained from the $A_\iy$ category $\cF(X)$ by taking the ($\fm^1$) cohomology groups of the hom spaces in it.
The key fact is that if $L_1,L_2$ are Hamiltonian isotopic then $L_1,L_2$ will define isomorphic objects in $H\cF(X)$ (after given the additional data including bounding cochains).

It is clear that we need only the weaker hypothesis because Thomas and Yau use in their proof only the Floer cohomology group $HF^*(L_1,L_2),$ the hom space in $H\cF(X).$
But when Thomas and Yan wrote their paper, the theory of Fukaya categories was much less developed at that time, so the statement was presumably better-sounding with Hamiltonian isotopies. 
In this paper we improve Thomas--Yau's theorem by using the more developed theory of Fukaya categories as follows.

{\bf(ii)}
Theorem \ref{01} is stronger than the original theorem in the respect that $L_1,L_2$ need not have (even after shift) the same phase. This would also follow from the following statement: two Lagrangians underlying the same isomorphism class of objects in $H\cF(X)$ should have the same homology class in $X;$ although we shall not discuss this in the present paper. 
The fact that $L_1,L_2$ have the same shift up to shift will be proved in Corollary \ref{7} (before giving the full proof of the theorem), for which we shall use some simple fact about Maslov indices and some nonvanishing results for Floer cohomology groups (which we recall in \S\ref{sec3}). 

Also we do not suppose that either $L_1$ or $L_2$ itself underlies an object of $H\cF(X)$
nor do we suppose even that it is cleanly immersed as Fukaya \cite{Fuk2} does.
But we do suppose that we can perturb $L_1,L_2$ both to generically immersed Lagrangians
which underlie objects of $H\cF(X).$
This is how we deal with badly immersed Lagrangians in the Floer-theory context in this paper.

{\bf(iii)}
The idea of the proof of Theorem \ref{01} is the same as that of Thomas and Yau. But we include the modification by Imagi, Joyce and Oliveira dos Santos \cite[Theorem 4.7]{IJO} because we can unfortunately not justify the original Morse-theory argument \cite[Theorem 4.3]{TY}.
\end{rem}

Another natural generalization is to
minimal Lagrangians in K\"ahler--Einstein manifolds
or more generally to $J$-minimal Lagrangians introduced by Lotay and Pacini \cite{LP,LP2}.
At the moment we state the result only in the K\"ahler--Einstein case.
For the more complete statement see Corollary \ref{2} (i).
\begin{thm}\l{02}
Let $X$ be a $\Z_k$-graded K\"ahler--Einstein manifold of complex dimension $n\le k;$
and $L_1,L_2\sb X$ two closed irreducibly-immersed $\Z_k$-graded minimal Lagrangians with the same grading on $L_1\cap L_2.$
Let $\b_1,\b_2\in H\cF(X)$ be two objects supported respectively near $L_1,L_2$ and such that:
\ea
\l{11}
&\text{either $\b_1\cong\b_2$ or $k\ge 2n-1$ and $\b_1\cong\b_2$ up to shift; and}\\
\l{12}
&HF^i(\b_1,\b_1)\cong HF^i(\b_2,\b_2)\ne0\text{ for } i=0,n.
\ea
Then $L_1=L_2\sb X.$
\end{thm}
\begin{rem}
The more precise meaning of $L_1,L_2$ having the same grading on $L_1\cap L_2$ is given in Definition \ref{dfn gr} below.
We make a more geometric definition of gradings than Seidel \cite{Sd1} does.
\end{rem}

In the circumstances of Theorem \ref{02},
if $c_1(X)>0$ there may be many automorphisms of $X$
and many zero objects of $H\cF(X)$ especially in the toric case as Fukaya, Oh, Ohta and Ono \cite{FOOOt} study.
% there is a certain potential function whose critical points correspond exactly to nonzero objects $H\cF(X).$
So we can hardly expect a general uniqueness statement to hold.
But for instance, to the {\it Clifford torus} 
\e\l{CT}T^n:=\{[z_0,\cdots,z_n]\in\CP^n:|z_0|=\cdots=|z_n|\}\e
we can apply Theorem \ref{02} with the latter alternative of \eq{11}.
We can in fact take $k=2n+2$ as we show in Example \ref{ExCliff} below. 
Also according to Fukaya, Oh, Ohta and Ono \cite{FOOOt}
there are $(n+1)$ objects $\b\in H\cF(X)$ supported on $T^n$
with $HF^*(\b,\b)\cong H^*(T^n).$ 
So \eq{12} holds and we have:
\begin{cor}\l{3}
Let $L\sb\CP^n$ be a closed irreducibly-immersed $\Z_{2n+2}$-graded minimal Lagrangian such that $L,T^n$ have the same grading on $L\cap T^n.$
Let $H\cF(X)$ have an object which is supported near $L$ and isomorphic to one of the $(n+1)$ objects $\b$'s above $($all supported on $T^n).$
Then $L=T^n.$
\end{cor}


On the other hand, if $c_1(X)\le0$ the nonzero condition \eq{12} holds automatically in certain circumstances.
For instance, if either $L_1$ or $L_2$ is embedded then \eq{12} holds automatically by
Fukaya, Oh, Ohta and Ono's theorem \cite[Theorem E]{FOOO}.
If $c_1(X)=0$ there is a stronger version of it which we use for Theorem \ref{01}. 

If $c_1(X)<0$ Theorem \ref{12} extends to $J$-minimal Lagrangians.
They are locally unique, so the result is interesting only in the global context.

For the proofs we follow in outline that by Thomas and Yau \cite[Theorem 4.3]{TY}.
In every case the heart of the proof is to
make certain Hamiltonian perturbations of the two Lagrangians.
The Floer-theory condition implies that the two perturbed Lagrangians have at least one intersection point of index $0$ or $n$
at which we do by analysis a sort of unique continuation.

But for the Hamiltonian perturbation process we use the method by
%we follow in outline the original proof by Thomas and Yau \cite[Theorem 4.3]{TY}
% their original Morse-theory method \cite[Theorem 4.3]{TY} looks difficult to justify.
Imagi, Joyce and Oliveira dos Santos \cite[Theorem 4.7]{IJO} as mentioned above.
The modified method works only in the real analytic category.
This causes no problem in the K\"ahler--Einstein case because
in that case everything is automatically analytic by elliptic regularity.
On the other hand, in Theorem \ref{01} and in the $J$-minimal version of Theorem \ref{02}
we work in the $C^\iy$ category
and the real analyticity condition is achieved by another Hamiltonian-perturbation process.
%For this we use the deformation theories by McLean \cite{ML} and Lotay--Pacini \cite{LP}
%respectively for special Lagrangians and for $J$-minimal Lagrangians.

We begin in \S\ref{sec2} with our geometric treatment of graded Lagrangians.
In \S\ref{sec2.5} we recall the relevant facts about Maslov forms in Lotay--Pacini's sense.
In \S\ref{sec3} we give a minimum account of Floer theory which we use in this paper. 
In \S\ref{sec4} we carry out the key step of making Hamiltonian perturbations.
In \S\ref{sec5} we state and prove all our results.
They are essentially at the chain level because \S\ref{sec4} is independent of Floer theory.


%In \S\ref{sec3} we give the minimum account of Floer theory which we use in 
%Apart from the nonvanishing theorems in \S\ref{sec3}
%we need only the standard fact that there exist Fukaya categories and Floer cohomology groups
%are invariant under Hamiltonian deformations.
%This is treated by Akaho and Joyce \cite{AJ},
%Fukaya \cite{Fuk2} and Fukaya, Oh, Ohta and Ono \cite{FOOO}
%although it will take some more time for everything to be written up.
%But , in certain circumstances we can use the $HF^*$ nonvanishing theorems in \S\ref{sec3} to deduce
%such nicer statements as Theorem \ref{01}.
%which is still an active research area including the works by
%Akaho and Joyce \cite{AJ}, Fukaya \cite{Fuk1,Fuk2} and Fukaya, Oh, Ohta and Ono \cite{FOOO,FOOOt,FOOO-Z}.
%to be given sooner or later a more complete treatment.
%But this does not mean that we shall need after all a big change for the results in this paper.

Finally we remark that Thomas--Yau's uniqueness theorem is just a bit of their whole proposal \cite{TY} which is now improved by Joyce \cite{Jc}.

{\it Acknowledgements.}
The author would like to thank Mohammed Abouzaid and Kenji Fukaya for helpful conversations.
This work was supported financially by the National Natural Science Foundation of China (11950410501).


\section{Geometric Gradings}\l{sec2}
We make a geometric definition of graded Lagrangians:
\begin{dfn}\l{dfn gr}
Let $(X,\om)$ a symplectic manifold, $J$ an $\om$-compatible almost-complex structure on $X$
and $K_X$ the canonical bundle over $(X,J).$
Let $k$ be either a positive integer or infinity, write $\Z_k:=\Z/k\Z$ for $k$ finite and write $\Z_\iy:=\Z.$
By a {\it $\Z_k$-grading} of $(X,\om,J)$ we mean for $k<\iy$ the pair of a complex line bundle $K_X^{2/k}$ and a bundle isomorphism $(K_X^{2/k})^k\cong K_X^2,$
and for $k=\iy$ a nowhere vanishing section $\Ps$ of $K_X^2.$ 
Suppose given such a $\Z_k$-grading of $X$ and a Hermitian metric $g$ on $(X,J).$
Then for every $x\in X$ and for every Lagrangian plane $\ell\sb T_xX$
there is a $g$-orthogonal decomposition $T_xX=\ell\op J\ell$ and accordingly
a canonical $g$-unit section of $K_X^2$ which we denote by $\Om_\ell^2;$
that is, if $e_1,\cdots,e_n\in\ell$ are a $g$-orthonormal basis then 
\[\Om_\ell^2:=\pm e^1\w\cdots\w e^n\w Je^1\w\cdots\w Je^n.\]
There is also a canonical $k$-fold cover $\Lag^kTX$ of the Lagrangian Grassmannian $\Lag TX$ defined for $k<\iy$ by
\[\Lag^kTX:=\{(\ell,\al)\in\Lag TX\t_X K_X^{2/k}:\Om_\ell^2=\al^k\}\]
and for $k=\iy$ by $\Lag^\iy TX:=\{(\ell,\ph)\in\Lag TX\t_X \R:\Om_\ell^2=e^{i\ph}\Ps\}.$ 
By a $\Z_k$-grading of an immersed Lagrangian $L\sb (X,\om)$
we mean a lift to $\Lag^kTX$ of the tangent-space map $L\to \Lag TX,$
that is, for $k<\iy$  a section $\al:L\to K_X^{2/k}$ with $\al^k=\Om_L^2$
and for $k=\iy$ a section $\ph:L\to S^1$ with $e^{i\ph}\Ps=\Om_L^2.$
It is unique up to $\Z_k$-shifts,
where the shift $[1]$ is defined for $k<\iy$ by $(L,\al)[1]:=(L,e^{2\pi i/k}\al)$
and for $k=\iy$ by $(L,\ph)[1]:=(L,\ph+2\pi).$
\end{dfn}
\begin{rem}
The advantage of this definition is that given two Lagrangians $L_i$ $(i=1,2)$ graded by $\al_i$ or $\ph_i$ as above,
we can compare the two gradings $\al_i$ or $\ph_i.$
In particular, it makes sense to say that they are equal or not.
\end{rem}

We make another definition in the circumstances of Definition \ref{dfn gr}:
\begin{dfn}
For $i=1,2$ let $L_i\sb(X,\om)$ be a Lagrangian with a $\Z_k$-grading $\al_i:L\to K_X^{2/k},$
$k<\iy.$
Suppose $L_1,L_2$ intersect transversely at a point $x\in L_1\cap L_2.$
We define the Maslov index $\mu_{L_1,L_2}(x)\in\Z_k.$
Define $e^{i\ta}\in S^1$ by $\al_2=e^{i\ta}\al_1\in K_X^{2/k}$ over $x.$
Take an isomorphism $(T_xX,J|_x,g|_x)\cong\C^n$
which maps $T_xL_1$ to $\R^n\sb\C^n$ and $T_xL_2$ to
$\{(e^{i\th_1}x_1,\cdots,e^{i\th_n}x_n)\in\C^n:x_1,\cdots,x_n\in\R\}.$
for some $\th_1,\cdots,\th_n\in(0,\pi).$
These $\th_1,\cdots,\th_n$ are unique up to order
and we have
$\Om_{L_2}^2=\pm e^{-2i(\th_1+\cdots+\th_n)}\Om_{L_1}^2.$
So we can define
\begin{equation}
\l{mu}\mu_{L_1,L_2}(x):=\frac{1}{2\pi}(2\th_1+\cdots+2\th_n+k\ta) \in\Z \text{ modulo }k\Z.
\end{equation}
If $k=\iy$ and if for $i=1,2$ the $L_i$ is $\Z$-graded by $\ph_i:L_i\to \R$ then
\e\l{mu2}\mu_{L_1,L_2}(x):=\frac{1}{2\pi}(2\th_1+\cdots+2\th_n-\ph_1(x)+\ph_2(x)) \in\Z.\e
\end{dfn}

We recall now the notion of {\it special} Lagrangians: 
\begin{dfn}\l{dfn slag}
Let $(X,\om)$ a symplectic manifold, $J$ an $\om$-compatible almost-complex structure on $X$
and $\Ps$ a $\Z$-grading of $(X,\om,J).$
We say that an immersed Lagrangian in $(X,\om,J,\Ps)$ is {\it special} if it is $\Z$-graded by a constant.
\end{dfn}
\begin{rem}
In the symplectic context including this definition and \S\ref{sec3} we do not need $J$ to be integrable.
But in the more geometric context we suppose $J$ integrable
and Harvey--Lawson's theorem applies as in Theorem \ref{HL}.
\end{rem}
%More generally, let $L,L'\sb X$ are two Lagrangians with $\Z$-gradings $\ph:L\to \R,$
%$\ph':L'\to \R$ respectively as above
%Suppose also that $L,L'$ intersect each other transversely at a point $x$
%and let $\th_1,\cdots,\th_n\in(0,\pi)$ be as in \eq{mu} which are well defined up to order.
%Then the Maslov index of $(L,L')$ at $x$ may be defined by
We prove a lemma which we use in Corollary \ref{7}:
\begin{lem}\l{6.9}
Let $(X,\om,J,\Ps)$ be such as in Definition \ref{dfn slag}
and $L_1,L_2\sb X$ two mutually-transverse special Lagrangian submanifolds which have not even after shifts the same grading.
Then there exists $i\in\Z$ such that the Maslov indices of $L_1,L_2$ all belong to $[i,i+n-1].$
This will still hold under Lagrangian perturbations of $L_1,L_2.$
\end{lem}
\begin{proof}
Let $\ph_1,\ph_2\in\R$ be gradings of $L_1,L_2.$
Then $\ph-\ph'\notin\Z$ and this condition is preserved under Lagrangian perturbations of $L_1,L_2.$ 
Let $x\in L_1\cap L_2$ and let $\th_1,\cdots,\th_n\in(0,\pi)$ be as in \eq{mu}.
Then by \eq{mu2} we have $\mu_{L_1,L_2}(x)\in[i,i+n-1]$ for some $i\in\Z.$
This $i$ is independent of $x$ which completes the proof. 
\end{proof}

We prove another lemma which we use in \S\ref{sec5}:
\begin{lem}\l{P0}
Let $(X,\om)$ be a symplectic manifold of real dimension $2n,$ $\Z_k$-graded with $k\ge n,$
equipped with a compatible almost complex-structure $J$ and equipped with a Hermitian metric $g.$
Let $L_1,L_2\sb X$ be two closed immersed $\Z_k$-graded Lagrangians with the same grading on $L_1\cap L_2.$
Denote by $S\sb L_1\cap L_2$ the set of points at which $L_1,L_2$ have at least one common tangent space.
Then for every open neighbourhood $U\sb X$ of $S$
there exists $\ep>0$ such that every Hamiltonian $\ep$-perturbation of $L_2$ that intersects $L_1$ generically $($that is, only at transverse double points$)$
has no intersection point with $L_1\cap U$ of index $0$ or $n$ modulo $k\Z.$
\end{lem}
\begin{proof}
If this fails there are an open neighbourhood $U\sb X$ of $S,$
a sequence of Hamiltonian perturbations $L_2^i$ of $L_2$ all transverse to $L_1$ and tending smoothly to $L_2,$
and a sequence of intersection points $x^i\in L_1\cap L_2^i\-U$ of index $0$ or $n.$
By hypothesis $L_1\cap L_2^i\-U$ is compact, so there is a subsequence of $x^i$ tending to some point $x\in L_1\cap L_2\-U.$
Suppose now that $k<\iy;$
the other case $k=\iy$ may be treated in the same manner.
Give $L_2^i$ a grading $\al^i$ by the obvious homotopy,
and write $\al^i=\exp(\sqrt{-1}\ta^i)\al$ over $x^i.$
By hypothesis $L_1,L_2$ have the same grading at $x^i$ so we may suppose that for each $i$ we have $|\ta^i|<2\pi/k$ and the sequence $\ta^i$ tends to $0$ as $i\to\iy.$
Denote by $\th_1^i,\cdots,\th_n^i\in(0,\pi)$ the $n$ angles between the two tangent spaces at $x^i.$
Then
$2\th_1^i+\cdots+2\th_n^i+k\ta^i=0\text{ or }n\pi\text{ modulo }2k\pi\Z.$
In fact, since $k\ge n$ it follows that this holds without modulo $2k\pi\Z.$
So the limits of the $n$ angles sum up to $0$ or $n\pi;$
that is, at the limit point $x$ there is a common tangent space to $L_1,L_2.$
But this contradicts the definition of $S.$
\end{proof}



\section{Maslov Forms}\l{sec2.5}
Following Lotay and Pacini \cite{LP2} we make:
\begin{dfn}\l{Mdfn}
Let $(X,\om,J)$ be a K\"ahler manifold of complex dimension $n$
and $K_X$ the canonical bundle over $(X,J).$
Let $L\sb (X,\om)$ be a Lagrangian submanifold, and $g$ a Hermitian metric on $(X,J)$
which need not be the K\"ahler metric of $(X,\om,J)$ nor even K\"ahler.
As in Definition \ref{dfn gr} take the canonical $g$-unit section $\Om_L^2$ of $K_X^2$ over $L.$
Denote by $\nb$ the Chern connection on $(X,J,g)$
and by $A$ the real $1$-form on $K_X^2$ over $L$ defined by
$\nb\Om_L^2=2iA\ot\Om_L^2.$
We call $A$ the {\it Maslov form} on $L$ relative to $g$
and say that $L$ is {\it $g$-Maslov-zero} if $A\=0.$
\end{dfn}
Here it is not essential that $J$ is integrable. Many results by Lotay and Pacini \cite{LP,LP2} hold with $J$ an $\om$-compatible almost-complex structure
and with $\nb$ a connection such that $\nb J=\nb g=0.$
But it is convenient for us to suppose $J$ integrable. We can then take the Chern connection; in \S\ref{sec4} we can use the underlying real analytic structure of $(X,J)$
which will be convenient for us to state our results; and there are as in Theorem \ref{HamPert1} nice deformation theory results for special Lagrangians and $J$-minimal Lagrangians which we use in Corollary \ref{3.1}.

As Lotay and Pacini \cite[Theorem 2.4]{LP} prove, if $g$ is K\"ahler the $g$-Maslov-zero condition and the $J$-minimal condition are equivalent.
If also $L$ is $g$-Lagrangian---that is, $g(Jv,v')=0$ for any
$v,v'\in TL$---then $L$ is $J$-minimal if and only if $L$ is $g$-minimal in the ordinary sense.

We work with $g$-Maslov-zero Lagrangians rather than with $J$-minimal Lagrangians because
for special Lagrangians we can take $g$ to be merely {\it conformally} K\"ahler as we recall now.
Let $(X,\om,J)$ a K\"ahler manifold and $\Ps$ a $J$-holomorphic volume form on $X.$
Joyce \cite{Jb} and other authors call $(X,\om,J,\Ps)$ an almost Calabi--Yau manifold but this does not mean that $J$ is merely an almost complex structure;
$J$ is integrable. It means that $\om$ is Andnot necessarily Ricci-flat.
%Following Joyce \cite{J5} we call an {\it almost Calabi--Yau} manifold a K\"ahler manifold $(X,\om,J)$ with a holomorphic volume form $\Om_X.$
Denote by $g$ the conformally K\"ahler metric on $(X,J)$ associated with $\ps^2\om$ where
$\ps:X\to\R^+$ is defined by
\e\l{ps}\ps^{2n}\om^n/n!=(-1)^{n(n-1)/2}(i/2)^n\Ps\w\ov\Ps.\e
Then for every Lagrangian submanifold $L\sb X$ the canonical $g$-unit section $\Om_L^2$
is of the form $\al\Ps^2|_L$ where $\al:L\to S^1$ is a smooth function.
This $L$ has zero Maslov $1$-form if and only if $\al$ is constant, that is, if and only if $L$ is special.
In these circumstances Harvey--Lawson's theorem \cite{HL} may be stated as follows:
\begin{thm}\l{HL}
If $L$ is special then $L$ is orientable and {\it calibrated} by $\sqrt\al\Ps$ in Harvey--Lawson's sense
where the choice of $\sqrt\al$ corresponds to the orientation of $L.$
Furthermore, $L$ is area-minimizing with respect to $g$ and $J$-minimal in $(X,J,g).$
\end{thm}

Lotay and Pacini \cite[Proposition 4.5]{LP2} prove a formula for the Maslov $1$-form $A$ and the $J$-mean-curvature vector $H_J$:
\e\l{AH}-JA=(H_J+T_J)\lrcorner g\e
where $T_J$ comes from the torsion of the Chern connection $\nb.$
If $L$ is special then by definition we have $A=0$ and $H_J$ is the ordinary mean curvature which is also zero. So by \eq{AH} we have $T_J=0$ too.
Conversely, if $L$ is merely $J$-minimal then $T_J$ may be nonzero.
This justifies us working with $g$-Maslov-zero Lagrangians rather than $J$-minimal Lagrangians.

%\begin{thm}\l{LPthm}
%In the circumstances above, suppose $g$ K\"ahler then
%$L$ is $g$-Maslov-zero if and only if $L$ is $J$-minimal.
%And {\bf(ii)} if also $L$ is Lagrangian relative to $g$---that is, $g(Jv,v')=$ for any
%$v,v'\in TL$---then $L$ is $J$-minimal if and only if $L$ is $g$-minimal in the ordinary sense.
%\end{thm}
We compute now the {\it relative first Chern class} $2c_1(X,L)\in H^2(X,L;\Z)$ for $g$-Maslov-zero Lagrangians $L\sb X.$
We recall therefore:
\begin{lem}
Let $(X,\om)$ be a symplectic manifold, $J$ an $\om$-compatible complex structure,
$K_X$ the canonical bundle over $(X,J)$
and $L\sb(X,\om)$ a Lagrangian submanifold.
Then the relative de-Rham class $2c_1(X,L)\in H^2(X,L;\R),$ which is integral, may be represented by the curvature $2$-form of any connection on $K_X^2$
that is flat over $L.$
\end{lem}
So by Definition \ref{Mdfn} we have: 
\begin{cor}\l{Masl}
Let $(X,\om,J)$ be a K\"ahler manifold,
$g$ a Hermitian metric on $(X,J),$
and $L\sb(X,\om)$ a $g$-Maslov-zero Lagrangian submanifold.
Then $2c_1(X,L)$ may be represented by
the curvature $(1,1)$-form of $K_X^2$ relative to the Chern connection on $(X,J,g).$
\end{cor}
\begin{exam}\l{ExCliff}
Let $X=\CP^n,$ $\om$ the Fubini--Study form and $g$ the Fubini--Study metric.
This is K\"ahler--Einstein and the curvature $(1,1)$-form of $K_X$ relative to $g$ is $-(n+1)\om.$
The Clifford torus $T^n\sb \CP^n$ defined by \eq{CT} is a minimal Lagrangian.
It is $J$-minimal and $g$-Maslov-zero so we can apply to it Corollary \ref{Masl}.

We prove that $\CP^n$ and $T^n$ may both be $\Z_{2n+2}$-graded.
Since $H_1(\CP^n,\Z)=0$ it follows according to Seidel \cite[Lemma 2.6]{Sd1} that $\CP^n$ may be $\Z_{2n+2}$-graded because $2c_1(\CP^n)=-2(n+1)[\om].$
Also $T^n$ may be $\Z_k$-graded if and only if $2c_1(\CP^n,T^n)$ is divisible by $k.$
By Corollary \ref{Masl} we have $c_1(\CP^n,T^n)=-(n+1)[\om]\in H^2(\CP^n,T^n;\R).$
Computation shows that $H_2(\CP^n,T^n;\Z)$ is generated by the image of $H_2(\CP^n,\Z)$
and the $n$ discs $D_a\sb\CP^n,$ $a\in\{1,\cdots,n\},$ defined by $|z_a|\le 1$ and $z_b =1$ for every $b\ne a.$
We have
\[\om|_{D_a}=\frac{ni}{2\pi}\frac{dz_a\w d\bar z_a}{(n+|z_a|^2)^2}
\text{ and }\int_{D_a}\om=1.\]
So $[\om]\in H^2(\CP^n,T^n;\R)$ is integral, q.e.d.
\end{exam}

We turn now to the deformation theory for $g$-Maslov-zero Lagrangians.
We recall Lotay--Pacini's lemma \cite[Lemma 4.1]{LP}
which we use in \S\ref{sec4}:
\begin{lem}\l{LP1}
Let $(X,\om,J)$ be a K\"ahler manifold
and $g$ a Hermitian metric on $(X,J).$
Let $L\sb (X,\om,J)$ be a closed immersed $g$-Maslov-zero Lagrangian
and $NL$ the $g$-normal bundle to $L\sb X.$
For every sufficiently small $v\in C^\iy(NL)$ 
denote by $Av$ the $g$-Maslov $1$-form on the graph of $v$
embedded in $X$ by the $g$-exponential map,
and denote by $A'$ the linearization of $A$ at $0\in C^\iy(NL).$
Then for every $v\in C^\iy(NL)$ we have
\e\l{A'}A'v=dd^*\hat v+\sum_{i=1}^ng(T_\nb(v,e_i),e_i)+v\lrcorner F_\nb\e
where $\hat v:=-Jv\lrcorner g,$ $\{e_1,\cdots,e_n\}$ is a $g$-orthonormal frame on $TL,$
$\nb$ the Chern connection of $(X,J,g),$
$T_\nb$ its torsion tensor and $F_\nb$ its curvature $(1,1)$-form for $K_X.$
\end{lem}
\begin{proof}
This is a version of Lotay--Pacini's lemma \cite[Lemma 4.1]{LP}.
They suppose that $(X,J,g)$ is almost K\"ahler
%; that is, $J$ need not be integrable but the Hermitian $(1,1)$-form of $(X,J,g)$ is closed.
but again this is not essential; the result applies to every almost Hermitian manifold $(X,J,g)$
with connection $\nb$ such that $\nb J=\nb g=0.$
Our formula \eq{A'} is simpler than Lotay--Pacini's in two respects.
One is that their $J$-volume function $\r_J:L\to\R^+$
is identically equal to $1$
and the other is that their projection operator $\pi_L$ is $g$-orthogonal.
These both hold because $L$ is $g$-Lagrangian,
which completes the proof.
\end{proof}

We remark that we can compute the torsion term in the conformally K\"ahler case.
If $g$ is a conformally K\"ahler metric on $(X,J)$ associated with $\ps^2\om$
then $\nb-\d \log\ps^2\ot\id$ is the Levi-Civita connection for $\om.$ 
Hence by computation we see that for any vector fields $u,v$ on $X$ we have
\[T_\nb(u,v)=d\log\ps\llcorner(u\ot v+Ju\ot Jv-v\ot u-Jv\ot Ju).\]
So the torsion term on \eq{A'} is equal to $(n-1)d\log\ps$ and
\e\l{A3}A'v=d(\ps^{n-1}d^*\ps^{1-n}\hat v)+v\lrcorner F_\nb.\e
This has application to deformation theory;
that is, by \eq{A3} we can apply Lotay--Pacini's result \cite[Proposition 4.5]{LP}
with $\ps^{n-1}$ in place of $\r_J.$
So Lotay--Pacini's uniqueness and persistence theorem \cite[Theorem 5.2]{LP} extends to
$g$-Maslov-zero Lagrangians with $g$ conformally K\"ahler.
But the corresponding perturbations are just Lagrangian and not necessarily Hamiltonian as we want
in the Floer theory context.

On the other hand, we can make (domain) Hamiltonian perturbations of special Lagrangians
and $J$-minimal Lagangians.
The result for special Lagrangians goes back to McLean's theorem \cite{ML}
and that for $J$-minimal Lagrangians is proved by Lotay--Pacini \cite[Theorem 5.6]{LP}.
\begin{thm}\l{HamPert1}
{\bf(i)}
Let $(X,\om,J)$ be a K\"ahler manifold, $\Ps$ a $J$-holomorphic volume form on $X$
and $L\sb X$ a closed immersed special Lagrangian relative to $(J,\Ps).$ 
%and $g$ the conformally K\"ahler metric on $X$ determined by $\om,\Ps$ as in Definition \ref{}.
Then for every sufficiently small perturbation $(J',\Ps')$ of $(J,\Ps)$ as $\om$-compatible complex structures and holomorphic volume forms relative to them,
there exists a domain Hamiltonian perturbation of $L$ which is special with respect to $(J',\Ps').$
\\{\bf(ii)}
Let $(X,\om,J)$ be a K\"ahler manifold such that $-\om$ is the Ricci $(1,1)$-form of another K\"ahler metric $g$ on $(X,J),$
and  $L\sb X$ a closed immersed $J$-minimal Lagrangian. 
Then for every sufficiently small perturbation $(J',g')$ of $(J,g)$ as $\om$-compatible complex structures and K\"ahler metrics relative to them
whose Ricci $(1,1)$-forms are all equal to $-\om,$
there exists a domain Hamiltonian perturbation of $L$ which is $J$-minimal with respect to $(J',g').$
\end{thm}


%\begin{cor}\l{LP2}
%In the circumstances above, if $g$ is a conformally K\"ahler metric on $(X,J)$ associated with $\ps^2\om$ for
%some $\ps:X\to\R^+$ then for every $v\in C^\iy(NL)$ we have
%\e\l{A2}A'v=d\ps^{n-1}d^*\ps^{1-n}\hat v+v\lrcorner F_\nb.\e
%\end{cor}
%\begin{proof}
%Letting $\nb_\om$ be the Chern connection of $(X,\om,J)$ and considering the effect of the conformal change, we see that
%\e\l{A3}\nb=:\nb_{\ps^2\om}=\nb_\om+\d \log\ps^2\ot\id.\e
%Recall now that for any two vector fields $u,v$ on $X$ we have
%\e\l{A4}T_\nb'(u,v)=\nb_u v'-\nb_v u'-[u,v]\e
%where the prime symbol are the complexifications $u':=\hf(u-iJu)$ and so on.
%Now by \eq{A3} and \eq{A4} we have
%\[
%T_\nb'(u,v)=\d\log\ps^2\llcorner(u\ot v'-v\ot u').\]
%And computation shows
%\[
%T_\nb(u,v)=d\log\ps\llcorner(u\ot v+Ju\ot Jv-v\ot u-Jv\ot Ju).
%\]
%In particular, $\sum_{i=1}^ng(T_\nb(v,e_i),e_i)=(n-1)d\log\ps.$
%Using this and recalling that $d^*=-*d*$ on $1$-forms, we see that
%\[dd^*\hat v+\sum_{i=1}^ng(T_\nb(v,e_i),e_i)=d\ps^{n-1}d^*\ps^{1-n}\hat v\]
%So by \eq{A'} we have \eq{A2}.
%\end{proof}
%Finally we prove:
%\begin{thm}\l{LP3}
%Let $(X,\om,J)$ be a K\"ahler manifold
%and $g$ a conformally K\"ahler metric on $(X,J)$ conformal to the K\"ahler metric of $(X,\om,J).$
%Suppose that the curvature $(1,1)$-form of $K_X$ relative to the Chern connection $\nb$
%on $(X,J,g)$ is positive definite.
%Then every closed immersed $g$-Maslov-zero Lagrangian in $(X,\om,J)$
%is locally unique and persists under perturbations of $(\om,J,g).$
%\end{thm}
%\begin{proof}
%Define a positive definite tensor $M:C^\iy(T^*L)\to C^\iy(NL)$ by $u=Mu\lrcorner F_\nb.$
%Fix $\al\in(0,1)$ and introduce for $l=0,1,2,\cdots$
%the Banach space $\cU^{l,\al}:=\{u\in C^{l,\al}(T^*L): du=0\}.$
%Define a linear operator $P:\cU^{2,\al}\to \cU^{0,\al}$ by $P:=A'M.$
%Then Lotay--Pacini's result \cite[Proposition 4.5]{LP} applies to our context with
%$\ps^{n-1}$ in place of $\r_J;$ that is,
%$P$ is elliptic and a Banach-space isomorphism.
%So by a standard application of the implicit function theorem we can prove the theorem. 
%\end{proof}



\section{$HF^*$ Nonvanishing Theorems}\l{sec3}
Following Fukaya, Oh, Ohta and Ono \cite{FOOO} given a unital commutative ring $R$ we define the Novikov ring
\e\l{La}\La=\left\{\sum_{i=0}^\iy a_iT^{\la_i}:a_i\in R[e,e^{-1}],\la_i\in\R\text{ for each }i\text{ and }\lim_{i\to\iy}\la_i=+\iy\right\}\e
where $T,e$ are two formal variables. They are related respectively to the areas and Maslov numbers of pseudoholomorphic curves.
We denote by $\La^0\sb\La$ the subring defined by the same formula \eq{La} as $\La$
but with $\la_i\ge0$ in place of $\la_i\in\R.$
We denote by $\La^+\sb\La^0$ the ideal defined by \eq{La} with $\la_i>0$ in place of $\la_i\in\R.$
We write $\px$ for $=$ modulo $\La^+.$

Fukaya categories are defined over $\La.$
If the ambient symplectic manifold is spherically nonnegative we can take $R=\Z$ as Fukaya, Oh, Ohta and Ono \cite{FOOO-Z} do.
But in general we have to do the pseudoholomorphic-curve counts over $\Q$ and we shall therefore suppose that $R$ contains $\Q.$
Otherwise, in our applications we need no condition on $R;$ that is, we work under the following:
\begin{hp}\l{0}
Let $(X,\om)$ be a $\Z_k$-graded symplectic manifold of real dimension $2n,$
either compact or $\om$-convex at infinity.
Let $R$ be a unital commutative ring such that if $X$ is not spherically nonnegative then $\Q\sb R.$
If $2\ne0$ in $R,$ we shall suppose that $X$ is given a background bundle for relative spin structures on Lagrangians in $X,$
and that the Lagrangians concerned are all given relative spin structures.
\end{hp}
Under this hypothesis there is a Fukaya category $\cF(X)$ and its cohomology category $H\cF(X).$
In this paper, following Akaho and Joyce \cite{AJ} we suppose that every object of $\cF(X)$ consists of
a closed generically-immersed Lagrangian $L\sb X$ and its additional structures such as bounding cochains.
More precisely, we should write the immersion map to $L,$ say $\Lh\to L,$ and `closed' means $\Lh$ being a  closed manifold.
We can include also a suitable class of local systems on $L$ and the results in this paper will hold for them but we shall omit this in our formal treatment.
We recall now the notion of bounding cochains now.

We recall first the homologically perturbed version of Floer complexes.
Given a closed manifold $\Lh$ and a generic Lagrangian immersion $\Lh\to L\sb X$ we write $L\cap L:=\Lh\t_X\Lh.$
This is the disjoint union of the diagonal $\Lh$ and self-intersection pairs.
It depends upon the choice of $\Lh\to X$ but is unique up to automorphisms of $\Lh.$
We write
$C_L^*:=H^*(L\cap L,R)[-\mu_{L,L}]$
where the diagonal has no degree shift and every self-intersection pair has degree equal to its Maslov index.
If $L$ is embedded, this $C_L^*$ is $\Z$-graded and supported in degrees $0,\cdots,n.$
If $L$ is $\Z_k$-graded, so is $C_L^*.$
If also $L$ has on its every self-intersection pair two gradings nearly equal,
then $C_L^*$ is supported in degrees $0,\cdots,n$ modulo $k\Z.$

We use the following version of $A_\iy$ algebra structure constructions.
There are on $C_L^0$ a unit and the intersection pairing so we can speak of unitality and cyclic symmetry as
Fukaya \cite{Fuk1} does.
 %Given on $(X,\om)$ a tame almost-complex structure $J$
%we introduce the discrete submonoid 
%$G_J\sb[0,\iy)\t\Z$
%generated by
%$(\be\cdot[\om],\be\cdot2c_1(X,L))$
%where $\be\in H_2(X,L;\Z)$ is a
%genus-zero stable $J$-holomorphic curve class.
\begin{thm}\l{A}
There is on $C_L^*\ot\La^0$ a cohomologically-unital cyclic
%gapped $\Z_k$-graded
curved $A_\iy$ algebra structure $(\fm^i)_{i=0}^\iy$ with $\fm^1\px0$ and $\fm^2\px\pm\w$ where $\w$ is the cup-product map.
If also $L$ is embedded, we can make this strictly unital. 
\end{thm}
\begin{rem}
The embedded case is proved by Fukaya \cite{Fuk1}.
This will extend to the immersed case if we give up having the strict unit and satisfy ourselves with
a cohomological unit.
To have the strict unit we ought to perturb in a certain manner the relevant pseudoholomorphic-curve moduli spaces.
But in the immersed case there would be moduli spaces of constant maps to
self-intersection points. If we perturbed these, the cyclic symmetry would fail.

The cyclic symmetry is used in Lemma \ref{v} below.
But also we give it another proof in which we use the open-closed map.
%from the Floer complex on $L$ to the ordinary complex on $X.$
\end{rem}

By a {\it bounding cochain} on $(C_L^*\ot\La^0;\fm^0,\fm^1,\fm^2,\cdots)$
we mean an element $\b\in C^1_L\ot\La^+$ with $\sum_{i=0}^\iy\fm^i(\b,\cdots,\b)=0.$
Given such a $\b$ there is a natural way of giving $C_L^*\ot\La^0$
a cohomologically-unital cyclic {\it ordinary} $A_\iy$ algebra structure
$(\fm^i_\b )_{i=0}^\iy$ with $\fm_\b^1\px0$ and $\fm_\b^2\px\pm\w.$

We denote by $1\in C_L^0$ the unit, which is the cohomological unit in Theorem \ref{A}.
We denote by $*1\in C_L^n$ the volume form supported on the diagonal,
which is if $2\ne0$ in $R$ the Poincar\'e dual to a point.
We denote also symbolically by
$(x,y):=\pm\int_{L\cap L}x\w y$
the intersection pairing on $C_L^*.$ The sign does not matter to us. 
We prove now:
\begin{lem}\l{v}
For every bounding cochain $\b\in C_L^1\ot\La^+$ we have:
\iz
\item[\bf(i)] if $L$ is embedded, we have $*1\notin\im\fm^1_\b;$ and
\item[\bf(ii)] if $\fm^1_\b *1=0$ then $[*1]\ne0$ in $H^n(C_L^*\ot\La^0,\fm_\b^1).$ 
\iz
\end{lem}
\begin{proof}
Suppose contrary to (i) that $*1=\fm^1_\b x$ for some $x\in C^*_L\ot\La.$
Then using the strict unit we find
\e\l{v0}
0=(\fm^1_\b 1,x)=\pm(1,\fm^1_\b x)=\pm(1,*1)\px\pm\int_{L\cap L}\fm^2_\b(1,*1)\px\pm\int_{L\cap L}*1=1.
\e
This is impossible which proves (i).

We give (ii) two proofs.
Firstly, if $*1=\fm^1_\b x$ for some $x\in C^*_L\ot\La$ then as in \eq{v0} we have
\e\l{v1}
0\px\pm\int_{L\cap L}\fm_\b^2(1,*1)
\e
but $\fm_\b^2(1,*1)$ need not be $*1.$
However $1$ is a cohomological unit and $*1$ an $\fm_\b^1$-cocycle, so
%there is $y\in C^*\ot\La$ such that
\e\l{v2}
\fm_\b^2(1,*1)=*1+\fm_\b^1y\px*1\px *1.%+dy
\e
%where $d$ is the differential on $C^*_L.$
Now by \eq{v1} and \eq{v2} we have
$0\px\pm\int_{L\cap L}*1\px\pm1\in R\-\{0\}.$
This is impossible and so completes the first proof.

In the other proof of (ii) we use neither the cyclic symmetry nor the unitality but the open-closed map constructed by Fukaya, Oh, Ohta and Ono \cite[\S3.8]{FOOO}.
This extends to generically immersed Lagrangians.
We use in particular the map
$\fp_\b:H^n(C_L^*\ot\La^0,\fm_\b^1)\to H^{2n}(X,R)\ot\La^0.$
Denote by $*_X1$ the volume form on $X$
and to distinguish it from that on $L$ write $*_L1:=*1\in C^n_L.$
Then by the extended version of Fukaya, Oh, Ohta and Ono's lemma \cite[Lemma 6.4.2]{FOOO}
we have $\fp_\b[*_L1]\px[*_X1].$
So $[*_L1]\ne0.$
\end{proof}

We prove now our version of Fukaya, Oh, Ohta and Ono's theorem \cite[Theorem E]{FOOO}.  
We recall therefore that the $A_\iy$ structure of Theorem \ref{A} is gapped;
that is, if we take on $X$ an $\om$-compatible  almost-complex structure $J$
we can write $\fm_\b^1$ as the sum of $\fm_{\b,\be}^1$ of degree $1-[\be]\cdot2c_1(X,L)$
where $\be\in H_2(X,L;\Z)$ is a genus-zero stable $J$-holomorphic curve class.
Also, given a bounding cochain $\b\in C_L^1\ot \La^+$ we define
$HF^*(\b,\b):=H^*(C_L^*\ot \La,\fm_\b^1)\ot\La.$
%and introduce the discrete submonoid 
%$G_J\sb[0,\iy)\t\Z$
%generated by
%$(\be\cdot[\om],\be\cdot2c_1(X,L))$
%where $\be\in H_2(X,L;\Z)$ is a
%genus-zero stable $J$-holomorphic curve class.

\begin{thm}\l{ne0}
Let Hypothesis \ref{0} hold and let $L\sb X$ be a closed generically-immersed Lagrangian.
Suppose also that:
\iz
\item[\bf(i)] $L$ is embedded and there is on $X$ an $\om$-compatible almost-complex structure $J$
such that for every $J$-holomorphic curve class $\be\in H_2(X,L;\Z)$
we have $\be\cdot2c_1(X,L)\le0;$ or
\item[\bf(ii)] $k\ge n+2$ and $L$ has on its every self-intersection pair the two gradings nearly equal.
\iz
Then for $i=0,n$ and for every object $\b\in H\cF(X)$ supported on $L$
we have $HF^i(\b,\b)\ne0.$
\end{thm}
\begin{proof}
The case (i) is due to Fukaya, Oh, Ohta and Ono \cite[Theorem E]{FOOO} but we recall the proof.
Since $L$ is embedded it follows that $C_L^*$ is $\Z$-graded
and supported in degrees $0,\cdots,n.$
Also, since $\be\cdot2c_1(X,L)\le0$ it follows that $\fm_\b^1$ has degree $\ge1.$
So under $\fm_\b^1$ nothing goes to $1$ and $*1$ goes to zero.
Since the cohomological unit is by definition an $\fm_\b^1$-cycle it follows now that $[1]\ne0$ in $H^0(C_L^*\ot\La^0,\fm_\b^1)$ and accordingly in $HF^0(\b,\b).$
On the other hand, we have proved that $*1$ is also an $\fm_\b^1$-cycle.
So by Lemma \ref{v} (i) we have $[*1]\ne0$ in 
$H^n(C_L^*\ot\La^0,\fm_\b^1)$ and accordingly in $HF^n(\b,\b).$

In the case (ii) we have $C_L^*$ merely $\Z_k$-graded.
But $k\ge n+2$ so $C_L^{-1}=C_L^{n+1}=0.$
Now $d_\b$ has degree $1$ modulo $k\Z$ so again
under $\fm_\b^1$ nothing goes to $1$
and $*1$ goes to zero.
So $[1]\ne0$ in $HF^0(\b,\b)$ and now by Lemma \ref{v} (ii) we have
$[*1]\ne0$ in $HF^n(\b,\b).$
This completes the proof.
\end{proof}


We apply Theorem \ref{ne0} to $g$-Maslov-zero Lagrangians.
We say that a closed immersed Lagrangian  $L\sb X$ {\it nearly} underlies $\cF(X)$ if there is an arbitrarily-small Hamiltonian deformation of $L$ which
underlies an object of $H\cF(X).$
Such an object is said to be {\it supported near $L$.} 
\begin{cor}\l{7.1}
Let Hypothesis \ref{0} hold with $X$ given an $\om$-compatible almost-complex structure $J.$
Let $g$ be another Hermitian metric on $(X,J)$ and
$L\sb (X,\om,J,g)$ a close $\Z_k$-graded immersed $g$-Maslov-zero Lagrangian.
Suppose that either:
{\bf(i)} $L$ is embedded; or
{\bf(ii)} $k\ge n+2$ and $L$ has on its every self-intersection pair the two gradings nearly equal.
Then for $i=0,n$ and for every object $\b\in H\cF(X)$ supported near $L$
we have $HF^i(\b,\b)\ne0.$ 
\end{cor}
\begin{proof}
In the case (i) it follows from Corollary \ref{Masl} that the condition (i) of Theorem \ref{ne0} holds.
This concerns only integral homology and cohomology classes and so is preserved under perturbations of $L.$
Taking one of them which underlies $\b$ we see from Theorem \ref{ne0} that $HF^i(\b,\b)\ne0$ as we want. 
Also in the case (ii) the condition (ii) of Theorem \ref{ne0} is preserved under perturbations of $L$
so the conclusion follows in the same way.
\end{proof}

For special Lagrangians we can say more:
\begin{cor}\l{7}
Let Hypothesis \ref{0} hold with $X$ given an $\om$-compatible almost-complex structure $J$ and a $J$-holomorphic volume form.
Then we have:
\iz
\item[\bf(i)]
For $i=0,n$ and for every object $\b\in H\cF(X)$ supported near a closed immersed special Lagrangian,
we have $HF^i(\b,\b)\ne0.$
\item[\bf(ii)]
Let $L_1,L_2\sb X$ be two closed immersed special Lagrangians
and let $H\cF(X)$ have two isomorphic objects $\b_1,\b_2$ supported respectively near $L_1,L_2.$
Then $L_1,L_2$ have up to shift the same grading.
\iz
\end{cor}
\begin{proof}
The part (i) follows immediately from Corollary \ref{7.1} with $k=\iy.$ 
If (ii) fails then by Lemma \ref{6.9} we have $HF^*(\b_1,\b_2)$ supported in degrees $i,\cdots,i+n-1$ for some $i\in\Z.$
So either $HF^0(\b_1,\b_2)=0$ or $HF^n(\b_1,\b_2)=0.$
But $\b_1\cong\b_2$ so either $HF^0(\b_1,\b_1)=0$ or $HF^n(\b_1,\b_1)=0.$
This however contradicts (i).
\end{proof}



\section{Hamiltonian Perturbation Theorem}\l{sec4}
In this section we prove the following theorem.
We work in the real analytic category
and for our applications in \S\ref{sec5} we can take the underlying real analytic structure of a complex manifold $(X,J).$
In fact any other real analytic structure will do if it makes everything analytic
but the statement will then be too long.
\begin{thm}\l{Ham}
{\bf(i)}
Let $(X,\om,J)$ be a real analytic K\"ahler manifold of complex dimension $n$
and $g$ a real analytic Hermitian metric on $(X,J).$
Let $L_1,L_2\sb X$ be two distinct closed irreducibly-immersed $g$-Maslov-zero Lagrangians
and let $S$ be as in Lemma \ref{P0}.
Then there exists a neighbourhood $U$ of $S\sb X$
and arbitrarily-small Hamiltonian deformations $L_1',L_2'$ of $L_1,L_2$ respectively
which intersect generically each other with no intersection point in $U$ of index $0$ or $n.$
This will still hold if $L_1,L_2$ are $g$-Maslov-zero only near $L_1\cap L_2.$
If also Hypothesis \ref{0} holds and $L_1,L_2$ nearly underlie $H\cF(X)$ we can make $L_1',L_2'$ underlie $H\cF(X).$
\\{\bf(ii)}
The assertion will hold even if $\om,g$ are not analytic but if there are
$\om$-compatible complex structures $J_t,$
$J_t$-Hermitian metrics $g_t$ and $g_t$-Maslov-zero Hamiltonian perturbations $L_{1t},L_{2t}$ of $L_1,L_2$
all parametrized smoothly by $t\in\R$ with $(J_0,g_0,L_{10},L_{20})=(J,g,L_1,L_2)$
and such that for $t\ne0$ the $\om,g_t$ are $J_t$-analytic.
\end{thm}
\begin{rem}
{\bf(i)}
The theorem is of local nature so in fact we need the $g$-Maslov-zero condition and the real analyticity condition only near $L_1\cap L_2.$
\\{\bf(ii)}
The index in $U$ may be defined without grading $L_1',L_2';$ that is, near $S$ we write $L_2'$
over $L_1'$ as the graph of an exact $1$-form $df$ and we take the Morse indices of $f.$
This definition agrees with that in \S\ref{sec2}
if $L_1,L_2$ are graded with the same grading on $L_1\cap L_2$ and if $L_1',L_2'$ are graded by the obvious homotopies.
%\\{\bf(iii)}
%
%Also we can make a version of the theorem in which $X$ is just a symplectic manifold,
%$g$ just a Riemannian metric on $X$ and $L,L'$ two ordinary minimal Lagrangians in $(X,g),$
%all analytic with respect to some real analytic structure on $X.$
%This is possible because the linearized operators for $J$-minimal Lagrangians and ordinary minimal Lagrangians have the same symbol.
\end{rem}

%We prove in \S\S\ref{40} and \ref{41} the key estimates which we use in \S\ref{42}.

%\subsection{Estimate by \L{}ojasiewicz}\l{40}
There are a few key estimates which we use for the proof of Theorem \ref{Ham}.
Firstly, according to \L{}ojasiewicz \cite{Lj}
for every function $f:\R^n\to\R$ with $f(0)=0$ and {\it analytic} near $0\in\R^n$
there exist two constants $c>0$ and $p\in(0,1)$ such that near $0\in\R^n$
we have $|f|^p\le c|df|.$
This implies readily that we have:
\begin{lem}\l{Loj}
%Let $X$ be a real analytic Riemannian manifold and $L,L'\sb X$ two analytic submanifolds of the same dimension.
%Let $S\sb L\cap L'$ be the set of points at which one of the tangent spaces to $L$ is equal to one of those to $L'.$
%Suppose $S$ compact.
In the circumstances of Theorem \ref{Ham} {\rm(i)}
take a Weinstein neighbourhood of $L_2\sb(X,\om)$ and write $L_1$ near $S$ as the graph over $L_2$ of a closed $1$-form $u.$
Then near $S$ we have
$|u|\ll |\nb u|\ll1;$
that is, for every $\ep>0$ there exists a neighbourhood of $S\sb L_2$ on which $|u|\le \ep|\nb u|\le\ep^2.$
Here $|\bl|,\nb$ are both computed with respect to the induced metric on $L_2.$
\end{lem}

%\subsection{Estimates for Maslov $1$-forms}\l{41}
We show next that the Maslov $1$-form operator may be well approximated by $dd^*$:
\begin{lem}\l{Hlem}
Let $(X,\om,J)$ be a K\"ahler manifold and let $g$ be a Hermitian metric on $(X,J).$
Let $L\sb X$ be a closed immersed $g$-Maslov-zero Lagrangian
and let $NL$ be the $g$-normal bundle to $L\sb X.$
Then there exists a constant $c>0$ such that for every sufficiently $C^1$-small $v\in C^\iy(NL)$
whose graph $($embedded in $X$ by the $g$-exponential map$)$ is $g$-Maslov-zero, we have at every point of $L$
\e\l{H0}|dd^*\hat v|\le c|v|_{C^1}(1+|v|_{C^2})\e
where $\hat v$ is the same as in Lemma \ref{LP1},
$|v|_{C^1}:=|v|+|\nb v|$
and $|v|_{C^2}:=|v|+|\nb v|+|\nb^2v|;$
the latter two are computed pointwise on $L$ with respect to $g.$
\end{lem}
\begin{proof}
Let $v\in C^\iy(NL)$ be sufficiently $C^1$-small
and $Av\in C^\iy(T^*L)$ the Maslov $1$-form of the graph of $v.$
The tangent space to the graph of $v$ involves the first derivative $\nb v$
so the canonical section of $K_X^2$ over the graph of $v$ is of the form $B(v,\nb v)$
where $B:NL\op(T^*L\ot NL)\to NL$ is a smooth function. 
Then $Av$ involves the differentiation of $B(v,\nb v)$ so 
\e\l{H0.5}Av=A_1(v,\nb v)+A_2(v,\nb v)\cdot\nb^2v\e
where $A_1:NL\op(T^*L\ot NL)\to NL$ and $A_2:NL\op(T^*L\ot NL)\to NL\ot TL\ot TL$ are some smooth functions.
Denote by $A'$ the linearization of $A$ at $0\in C^\iy(NL).$
Then by \eq{H0.5} there is a $v$-independent constant $c_1>0$ such that
\e\l{H1}|(A-A')v|\le c_1[|v|^2+|\nb v|^2+|\nb^2 v|(|v|+|\nb v|)]\le c_1|v|_{C^1}|v|_{C^2}\e
where the same constant $c_1$ will do for both the inequalities.
Now by Lemma \ref{LP1} $A'$ is, up to lower-order terms, equal to $dd^*\hat v;$
that is, there is a $v$-independent constant $c_2>0$ such that
\[|A'v-dd^* \hat v|\le  c_2|v|_{C^1}.\]
So by \eq{H1} we have $|(A-A')v|\le \max\{c_1,c_2\}|v|_{C^1}(1+|v|_{C^2}).$
This implies indeed that if $Av=0$ then \eq{H0} holds.
%On the other hand we have
%\[P'v=\d_1P_1(0,0)\cdot v
%+\d_2P_1(0,0)\cdot\nb v
%+P_2(0,0)\cdot\nb^2\]
%where $\d_1,\d_2$ denote the partial differentiations in the fibres of $NL,(T^*L\ot NL)$ respectively.
%\e\l{H2.5} \d_2P(0,0)\cdot\nb v+\nb^2v\cdot Q(0,0)=d(\r_J^{-1}d^*\r_J \hat v)\e
%********
%The computation by Lotay and Pacini \cite[Lemma 4.1]{LP} shows that
%the first and second order terms of $P'$ is $d(\r_J^{-1}d^*\r_J \hat v)$
%where $\r_J:L\to\R^+$ is the $J$-volume function;
%but all we need is that $P'$ is, up to lower-order terms, equal to $dd^*\hat v.$
%This combined with \eq{H1} implies \eq{H3}.
%********
%There is also a certain formula for $\d_1P(0,0).$
\end{proof}
% as Chen and Oh, Tsai--Wang do but we shall not need that. 
%Its linearization at $0\in C^\iy(NL)$ is given by the second-variation formula for the area functional,
%so that for every $v\in C^\iy(NL)$ we have
%\[d_0Pv=\nb^*\nb v-\sum_{i,a,b}R_{iaib}v_a-\sum_{i,j,a,b}A_{ija}A_{ijb}v_a\]
%where $i,j$ are indices for $TL$ and $a,b$ for $NL$ both relative to orthonormal frames.    
%If moreover $X$ is K\"ahlerian and $L$ is Lagrangian, it is well known that
%\[-{\rm Ric}_X(v,v)+{\rm Ric}_L(Jv,Jv).\]
%For applications of this see ???.
We give a corollary to the lemma above:
\begin{cor}\l{Hcor}
In the circumstances of Lemma \ref{Hlem}
take a Weinstein neighbourhood of $L\sb X.$
Then there exists a constant $c>0$ such that for every sufficiently $C^1$-small $1$-form $u\in C^\iy(T^*L)$ whose graph is also minimal, we have at every point of $L$
\e\l{H5}|dd^* u|\le c|u|_{C^1}(1+|u|_{C^2}).\e
\end{cor}
\begin{proof}
The two embeddings, one given by the exponential map and the other given by the Weinstein neighbourhood, are mutually related by a diffeomorphism $F$ between neighbourhoods of the zero sections of $T^*L,NL$ respectively which induces the identity on the zero section $L,$ and over every point of $L,$ a diffeomorphism of the two fibres.
So there is a smooth function $F_1:T^*L\to NL\ot TL$
such that for every sufficiently $C^1$-small $u\in C^\iy(T^*L)$ we have
$Fu=F_1u\cdot u.$
In particular there is a $u$-independent constant $c>0$ such that
\e\l{H6} |Fu|_{C^1}\le |u|_{C^1}\text{ and }|Fu|_{C^2}\le c|u|_{C^2}.\e
The function $Gu:=-JFu\lrcorner g$ is also of the same form; that is,
there is a smooth function $G_1:T^*L\to T^*L\ot TL$ such that
$Gu=G_1u\cdot u.$
This implies in turn that we have, making $c$ large enough,
\e\l{H7}
|dd^*Gu|\le c|dd^*u|+c|u|_{C^1}(1+|u|_{C^2})
\e
Now by \eq{H0} we have, making $c$ large enough,
\[ |dd^* Gu|\le c|Fu|_{C^1}(1+|Fu|_{C^2}).\]
So by \eq{H7} we have \eq{H5} with $c$ large enough.
\end{proof}

%\subsection{Proof of Theorem \ref{Ham}}\l{42}
Now we prove Theorem \ref{Ham} in three steps:
\begin{proof}[Step 1: proof of the first part of {\rm(i)}]
In this case we take $L_1'=L_1$ and perturb only $L_2.$
We take a Weinstein neighbourhood of $L_2\sb (X,\om)$
and define $L_2'$ as the graph of a certain exact $1$-form $df$ on $L_2$ with $f:L_2\to\R$ taken as follows.
Since $L_1,L_2$ are both connected and analytic with $L_1\ne L_2$ as subsets of $X$ it follows that $L_2\-L_1$ is dense in $L_2.$
So there is a Morse function $f:L_2\to\R$ with $df\ne0$ on $L_1\cap L_2.$
Since $S\sb L_1\cap L_2$ it follows then that
there is a constant $c>0$ such that we have near $S$
\e\l{f}|f|_{C^2}<c|df|.\e
This condition is preserved under rescalings and smooth perturbations of $f$
so we can make $L_2'$ arbitrarily $C^\iy$ close to $L_2$ and transverse to $L_1.$

Near $S$ take a local component of $L_1$ if need be (in the strictly immersed case) and write it as the graph over $L_2$ of some closed $1$-form $u.$
This is possible by the definition of $S$ and in fact we have also $u=\nb u=0$ on $S.$
So by Lemma \ref{Loj} we have near $S$
\e\l{ph1}|u|\ll|\nb u|\ll1.\e
On the other hand, by Corollary \ref{Hcor} with $L_2$ in place of $L,$ there is a constant $c'>0,$ which depends upon $|u|_{C^2},$
such that we have near $S$
\e\l{ph2} |dd^*u|\le c' (|u|+|\nb u|)\le 2c'|\nb u|\ll1\e
where the second inequality follows from \eq{ph1}.
Applying to $d^*u$ the \L{}ojasiewicz estimate, we find a constant $\ep>0$ such that we have near $S$
\e\l{ph3} |d^*u|\le |dd^*u|^{1+\ep}\le (2c'|\nb u|)^{1+\ep}\ll|\nb u|\e
where the second inequality follows from \eq{ph2}.

We show now that we have on $L_1\cap L_2'$ near $S$
\e
\l{Hs}|d^*(u-df)|\le n^{-1/2}|\nb(u-df)|.
\e
On $L_1\cap L_2'$ near $S$ we have $u=df$ and so by \eq{f} 
\e\l{phf}|\nb(u-df)|\ge|\nb u|-|\nb df|\ge|\nb u|-c|df|=|\nb u|-c|u|\ge\frac{1}{2}|\nb u|\e
where the last inequality follows from \eq{ph1}.
On the other hand by \eq{ph3} we have on $L_1\cap L_2'$ near $S$
\e\l{phf2}
|d^*(u-df)|\le|d^*u|+c|df|
\le\frac{1}{4}n^{-1/2}|\nb u|+c|u|
\le\frac{1}{2}n^{-1/2}|\nb u|
\e
where the last inequality follows again from \eq{ph1}.
By \eq{phf} and \eq{phf2} we have \eq{Hs} as we want.

Finally we take a point $x\in L_1\cap L_2'$ near $S$
and compute its index,
which is the index of the nondegenerate symmetric bilinear form $\nb(u-df)$ on $T_xL_2.$
Diagonalize it by an orthonormal basis of $T_xL_2$ and denote by $\la_1,\cdots,\la_n$ the diagonal entries.
Then by \eq{Hs} we have
\[\ts|\la_1+\cdots+\la_n|\le n^{-1/2}(\la_1^2+\cdots+\la_n^2)^{1/2}\le\max\{|\la_1|,\cdots,|\la_n|\}.\]
But $\la_1,\cdots,\la_n$ are all nonzero so they have not all the same sign; that is, the index is neither $0$ nor $n$ as we want.
\end{proof}

\begin{proof}[Step 2: proof of the latter part of {\rm(i)}]
We make first a generic Hamiltonian perturbation of $L_2$ which underlies an object of $H\cF(X)$
and we define $L_2'$ as a further Hamiltonian perturbation of it.
As in step 1 we take a Weinstein neighbourhood of $L_2$ and define $L_2'$ as the graph of some exact $1$-form $dh_2$ over $L_2$
with $h_2\ne0$ on $S.$
Also as in step 1 write $L_1$ near $S$ as the graph of some real analytic closed $1$-form $u$ on $L_2$ near $S.$
Take a generic Hamiltonian perturbation $L_1'$ of $L_1$ which underlies $H\cF(X)$
and is close enough to $L_1$ to be written near $S$ as the graph of $u+dh_1$ over $L_2$ with
\e\l{ph5}|h_1|_{C^2}\ll |h_2|_{C^2}.\e
Since $dh_2\ne 0$ on $S$ it follows that there is a constant $c>0$ such that near $S$
we have $|h_2|_{C^2}<\frac{c}{4}|dh_2|.$
And then putting $f:=h_2-h_1$ we have
\e|h_2|_{C^2}<\frac{c}{4}(|df|+|dh_1|)\le\frac{c}{4}|df|+\frac{1}{2}|h_2|_{C^2}\l{ph6}\e
where the last inequality follows from \eq{ph5}. The whole estimate \eq{ph6} implies
$|h_2|_{C^2}<\frac{c}{2}|df|.$
Hence using again \eq{ph5} we find
\[|f|_{C^2}\le |h_2|_{C^2}+|h_1|_{C^2}<
\frac{c}{2}|df|+|h_2|_{C^2} <c|df|;\]
that is, the estimate \eq{f} holds as in step 1.
Following the subsequent estimates we see also that \eq{Hs} holds
now on $L_1'\cap L_2'$ near $S.$
This implies again that $L_1',L_2'$ have near $S$ no intersection point of index $0$ or $n.$
\end{proof}

\begin{proof}[Step 3: proof of Theorem \ref{Ham} {\rm (ii)}]
As in step 2 we make Hamiltonian perturbations $L_1',L_2'$ of $L_1,L_2$ both underlying $H\cF(X).$
The $L_2'$ is the graph of $dh_2$ over $L_2$ and
the $L_1'$ near $S$ is the graph of $u+dh_1$ over $L_2$ near $S.$
The difference $f:=h_2-h_1$ satisfies \eq{f} as in step 1.
We extend these to smooth families. By Theorem \ref{HamPert1} there are $g_t$-Maslov-zero Hamiltonian perturbations
$L_{1t},L_{2t}$ of $L_1,L_2.$
Extend the Weinstein neighbourhood of $L_2$ to those of $L_{2t}$
and write $L_2'$ as the graph of some $dh_{2t}$ over $L_{2t}.$
Denote by $S_t\sb L_{1t}\cap L_{2t}$ the set of intersection points at which $L_{1t},L_{2t}$ have at least one common tangent space.
Write $L_{1t}$ near $S_t$ as the graph of a real analytic closed $1$-form $u_t$ over $L_{2t}$ near $S_t$
and write $L_1'$ near $S_t$ as the graph of $u_t+dh_{1t}$ over $L_{2t}$ near $S_t.$
Then with $t$ small enough the estimate \eq{f} holds for $f_t:=h_{2t}-h_{1t}$ in place of $f$
and we can follow the subsequent estimates in step 1.
We can do this in a $t$-independent neighbourhood $U$ of $S\sb X$ because $S_t$ tends to $S.$
So $L_1',L_2'$ have in $U$ no intersection point of index $0$ or $n.$
\end{proof}
%\begin{proof}[Step 3]
%Finally in the immersed case we take immersions $\Lh_i\to L_i\sb X$ for $i=1,2$
%and work in a Weinstein neighbourhood of $\Lh_2$
%because Weinstein's theorem applies to $\Lh_2\to X.$
%The pull-back of $S$ in $\Lh_1$ is compact and so has finitely many components.
%Its every component may be treated as above, which completes the proof.
%\end{proof}

\section{Conclusions}\l{sec5}
From Lemma \ref{P0} and Theorem \ref{Ham} we deduce:
\begin{thm}\l{P1}
Let $(X,\om,J)$ be a real analytic K\"ahler manifold of complex dimension $n,$ $\Z_k$-graded with $k\ge n$ and given 
another real analytic Hermitian metric $g.$
Let $L_1,L_2\sb X$ two distinct closed irreducibly-immersed $g$-Maslov-zero Lagrangians with the same grading on $L_1\cap L_2.$
Then there exist arbitrarily-small Hamiltonian deformations $L_1',L_2'$ of $L_1,L_2$
which intersect generically each other with no intersection point of index $0$ or $n.$
If also Hypothesis \ref{0} holds and $L_1,L_2$ nearly underlie objects of $H\cF(X)$ then we can make $L_1',L_2'$ underly objects of $H\cF(X).$
\end{thm}
\begin{proof}
Let $U$ be as in Theorem \ref{Ham} and corresponding to this $U$ let $\ep$ be as in Lemma \ref{P0}.
Then corresponding to this $\ep$ let $L_i'$ $(i=1,2)$ be such $\ep$-perturbations of $L_i$ as in Theorem \ref{Ham}.
Applying Lemma \ref{P0} to the intersection points outside $U$
and Theorem \ref{Ham} to those in $U$
we see that $L_1',L_2'$ are such as we want.
\end{proof}


We make now the $C^\iy$ version of Theorem \ref{P1}:
\begin{thm}\l{3.1}
Let $(X,\om,J)$ be a K\"ahler manifold of complex dimension $n$
which is either: $\Z$-graded by a $J$-holomorphic volume form $\Ps$ and given the conformally K\"ahler metric $g$ as in $\S\ref{sec2}$ by using $\eq{ps};$
or $\Z_k$-graded with $k\ge n$ and given another K\"ahler metric $g$
whose the Ricci $(1,1)$-form is $-\om.$ 
Let $L_1,L_2\sb (X,\om,J,g)$ be two distinct closed irreducibly-immersed $g$-Maslov-zero Lagrangians with the same grading on $L_1\cap L_2.$
Suppose that one of the following three conditions holds: {\bf(i)} $X$ is compact;  {\bf(ii)} $X$ is $\om$-convex, $\om$ is exact and $J$ is Stein; or
{\bf(iii)} either $L_1$ or $L_2$ is generically immersed.
Then the same conclusion holds as in Theorem \ref{P1} (both the two sentences in it).
\end{thm}
\begin{proof}
We reduce the problem to the real analytic case by a further perturbation process.
In the case (i) we use the fact that
for every compact K\"ahler manifold $(X,\om,J)$
there exists on $(X,J)$ a smooth family $(\om_t)_{t\in\R}$ of K\"ahler forms with $\om_0=\om$ and $\om_t,$ $t\ne0,$ analytic.
This seems well known but for clarity we give it a proof.
Denote by $g_\om$ the K\"ahler metric of $(X,\om,J)$
and take on $(X,J)$ a smooth family of {\it Riemannian} metrics $g_t'$
with $g_0'=g_\om$ and $g_t',$ $t\ne0,$ analytic.
Then $g_t:=\hf(g_t'+J^*g_t')$ defines on $(X,J)$ a smooth family of Hermitian metrics with
$g_0=g_\om$ and $g_t,$ $t\ne0,$ analytic.
Denote by $\d_t^*,\db_t^*$ the respective formal $g_t$-adjoints of $\d,\db$
and following Kodaira and Spencer \cite[\S6]{KS}
introduce the smooth family of elliptic operators
\[E_t:=\d\db\db_t^*\d_t^*+\db_t^*\d_t^*\d\db+\db_t^*\d\d_t^*\db+
\d_t^*\db\db_t^*\d+\d_t^*\d+\db_t^*\db.\]
Since $g_0$ is K\"ahler it follows readily that 
$E_0$ is apart from the last two terms just the squared Laplacian.
The last two terms are so added that the kernel of $E_t$ consists of closed forms.
Denote by $\om_t'$ the projection of $\om$ onto the kernel of $E_t$ in the space of $(1,1)$-forms.
Then by Kodaira--Spencer's result \cite[Proposition 8]{KS} the kernel of $E_t$ has dimension independent of $t.$
So by another result of them \cite[Theorem 5]{KS} $\om_t'$ is smooth with respect to $t.$
Putting $\om_t:=\hf(\om_t'+\ov\om_t')$ we see that $\om_t$ is such as we want.

In the case (ii) we write $\om=dJdp$ for some $p:X\to\R$ and perturb $p$ to a real analytic function.
In the case (iii) we take a Stein open set in $(X,J)$ on which $\om$ is exact and which contains the generically immersed Lagrangian, say $L_1.$
This is possible because every self-intersection point of $L_1$ is a transverse double point;
the Stein structure near it is given by taking the product of squared distances from the two components of the generically immersed Lagrangian, which is a plurisubharmonic function.  

In every case there are K\"ahler perturbations $\om_t$ of $\om$ and by Moser's theorem diffeomorphisms $\Ph_t:X\to X$
with $\om=\Ph_t^*\om_t.$
Then $\om$ is analytic with respect to $J_t:=\Ph_t^*J$
and for $i=1,2$ the Lagrangian $\Ph_t^*L_i\sb(X,\om)$ is a nearly-zero Maslov-form
relative to $g_t:=\Ph_t^*g.$
So by Theorem \ref{HamPert1} there is a $g_t$-Maslov-zero Hamiltonian perturbation $L_{it}$ of $\Ph_t^*L_i$
to which we can apply Theorem \ref{Ham} (ii).
Combining it with Lemma \ref{P0} as in the proof of Theorem \ref{P1} we see that Theorem \ref{3.1} holds.  
\end{proof}

%\begin{cor}\l{5}
%Let $(X,\om,J)$ be a K\"ahler manifold with $c_1(X)<0$
%and let $-\om$ be the Ricci $(1,1)$-form of another K\"ahler metric $g$ on $(X,J).$
%either:
%\iz
%\item[\bf(i)]
%there is on $(X,\om,J)$ a conformally rescaled metric $g$
%such that the curvature $(1,1)$-form of the canonical bundle $K_X$ relative to the Chern
%connection on $(X,J,g)$ is positive definite; or
%\item[\bf(ii)]
%$-\om$ is the Ricci $(1,1)$-form of another K\"ahler metric $g$ on $(X,J).$
%\iz
%\begin{rem}
%Lotay--Pacini's theorem \cite[Theorem 5.2]{LP} extends to the circumstances of (i)
%as we show in Theorem \ref{LP3}.
%For (ii) we use their theorem \cite[Theorem 5.6]{LP} as it is,
%in which case $g$ is K\"ahler so $L,L'$ are $J$-minimal.
%\end{rem}
Finally using the results of \S\ref{sec3} we get:
\begin{cor}\l{2}
{\bf(i)}
Let $(X,\om,J,g,L_1,L_2)$ be such as in Theorem \ref{P1} or Theorem \ref{3.1} except that $L_1,L_2\sb X$ may be equal.
Let Hypothesis \ref{0} hold and let $\b_1,\b_2\in H\cF(X)$ be two objects supported respectively near $L_1,L_2$ such that:
\e\l{21}
\text{either $\b_1\cong\b_2$ or $k\ge 2n-1$ and $\b_1\cong\b_2$ up to shift.}
\e
Suppose also that 
\e
\l{22}
HF^i(\b_1,\b_1)\cong HF^i(\b_2,\b_2)\ne0\text{ for }i=0,n.
\e
More generally, suppose as in Corollary \ref{7.1} that either: $c_1(X)\le0$ and one of $L_1,L_2$ is embedded;
or $k\ge n+2$ and one of $L_1,L_2$ has on its every self-intersection pair the two gradings nearly equal.
Then $L_1=L_2\sb X.$
\\{\bf(ii)}
Let $(X,\om,J)$ be a K\"ahler manifold equipped with a holomorphic volume form;
and $L_1,L_2\sb X$ two closed irreducibly-immersed special Lagrangians.
Suppose that either $\om$ is $J$-analytic or one of {\rm (i)--(iii)} of Theorem \ref{3.1} holds. 
Let Hypothesis \ref{0} hold and let $H\cF(X)$ have two isomorphic objects supported respectively near $L_1,L_2.$
Then $L_1=L_2\sb X.$
\end{cor}
\begin{proof}
Suppose contrary to the assertion that $L_1\ne L_2.$
Then perturbing $L_1,L_2$ as in Theorem \ref{P1} or Theorem \ref{3.1} 
we see that $HF^*(\b_1,\b_2)$ is supported in degrees $1,\cdots,n-1$ modulo $k\Z.$
But by hypothesis we have $k\ge n$ so $HF^0(\b_1,\b_2)=HF^n(\b_1,\b_2)=0.$ 
If also $\b_1\cong\b_2$ then $HF^0(\b_1,\b_1)=HF^n(\b_1,\b_1)=0$ which however contradicts \eq{22} or Corollary \ref{7.1}. This is the first case of \eq{21}.

In the second case we have $k\ge2n-1$ but $\b_1\cong\b_2$ only up to shift.
Then from Theorem \ref{P1} we see only that
$HF^*(\b_1,\b_1)$ is supported in degrees $i,\cdots,i+n-2$ for some $i\in\Z.$
But again we have $HF^0(\b_1,\b_1)\ne0$ and $HF^n(\b_1,\b_1)\ne0$ so after translating $[i,i+n-2]$ by $k\Z$
we may suppose $0\in[i,i+n-2]$ and $n\in[i+j,i+j+n-2]$ for some $j\in k\Z.$
Then $i\in[2-n,0]\cap[2-j,n-j]$ so $0\ge 2-j$ and $2-n\le n-j;$ that is, $j\in[2, 2n-2].$
So $0<j<2n-1\le k$ which however contradicts $j\in k\Z.$
This proves Corollary \ref{2} (i).

For Corollary \ref{2} (ii) we use Corollary \ref{7} (ii) to see that $L_1,L_2$ have up to shifts the same grading.
We can then apply Corollary \ref{2} (i) with $k=\iy$ which completes the proof.
\end{proof}








% %   \end{equation}
% %   The space $\Obs(N, \rho) $ is equipped with a canonical element called $0$.
% % %  the moduli space of holomorphic discs associated to a compatible almost complex structure $J$, which is conical at infinity, determines a function from $H^1(N)$ to $H^2(N)$. The zero-locus of this function is called the set of bounding cochains.
% % \item The triples $(N, \rho, b)$, where $b \in \Def(N, \rho)$ is an element mapping to $0$ in $\Obs(N, \rho)$ are the objects of an  $A_\infty$ category (called the wrapped Fukaya category). We call such a triple a \emph{brane supported on $N$}.
% %   For each representation $\rho$ of $\pi_1(Q)$ which is isomorphic to the tensor product of an abelian and a unitary representation (see Definition \ref{def:Lagrangian-brane}), there is an object of the Fukaya category whose Floer cohomology with the cotangent fibre is isomorphic to $\rho$, and which is supported on the graph of a closed $1$-form.


%   In the simpler context of exact Lagrangians, the first 




% the most general any of the variants of this category, the objects are immersed Lagrangians (with transverse self-intersections) equipped with auxiliary data. Its morphisms are 

% We say that an object 









% As mentioned earlier, unobstructed Lagrangians are those which support objects of the Fukaya category.  This is a Floer-theoretic notion, in the sense that it invariant under all symplectic diffeomorphisms, but that its formulation requires the choice of auxiliary data, the most substantial of which is the choice of an almost complex structure $J$ which is compatible with the symplectic form in the sense that the $2$-tensor $\omega(J \cdot, \cdot)$ is a metric. Because the cotangent bundle is not compact, it will be important for technical reasons to ensure that this almost complex be of contact type as discussed in Section \ref{FT}.
% \begin{rem}
% In our applications, we start with an integrable complex structure near the $0$-section, arising from an embedding of a Weinstein neighbourhood in a Calabi--Yau manifold. Such a complex structure may not have an integrable extension to the entire cotangent bundle, but we can always an extension to a compatible almost complex structure which is of contact type.
% \end{rem}





% \begin{enumerate}
% \item Associated to each closed relatively $\Spin$ immersed Maslov-$0$ Lagrangian $N$ of $T^* Q$ with transverse self-intersections, and each choice of integral local system $\rho$ on $N$, there are a pair of sets $\Def(N, \rho) $ and $\Obs(N, \rho)$, and a map
%   \begin{equation}
%     \Def(N, \rho) \to \Obs(N, \rho).
%   \end{equation}
%   The space $\Obs(N, \rho) $ is equipped with a canonical element called $0$.
% %  the moduli space of holomorphic discs associated to a compatible almost complex structure $J$, which is conical at infinity, determines a function from $H^1(N)$ to $H^2(N)$. The zero-locus of this function is called the set of bounding cochains.
% \item The triples $(N, \rho, b)$, where $b \in \Def(N, \rho)$ is an element mapping to $0$ in $\Obs(N, \rho)$ are the objects of an  $A_\infty$ category (called the wrapped Fukaya category). We call such a triple a \emph{brane supported on $N$}.
% \item  If $N'$ is Hamiltonian isotopic to $N$, then for each brane supported on $N$, there is a quasi-isomorphic brane supported on $N'$.
% \item The wrapped Fukaya category is generated by the cotangent fibre
% \item The self-Floer cohomology of the cotangent fibre is supported in non-positive degree, and is isomorphic in degree $0$ to the group ring of $\pi_1(Q)$.
%   \item For each representation $\rho$ of $\pi_1(Q)$ which is projectively unitary, there is a brane $\cQ_{\rho}$, supported on a the graph of a closed $1$-form on $Q$, with the property that the Floer cohomology $\cQ_{\rho}$ with a cotangent fibre is isomorphic to $\rho$.
% \end{enumerate}









% We consider a cotangent bundle $T^* Q$ as our ambient symplectic manifold, equipped with a compatible almost complex structure, which is conical at infinity, and the standard trivialisation of the complex line bundle given by the square of the $n$th (complex) exterior power of the tangent space, considered as an almost complex vector bundle. 

% We shall rely on the following facts:
% \begin{enumerate}
% \item Associated to each closed relatively $\Spin$ immersed Maslov-$0$ Lagrangian $N$ of $T^* Q$ with transverse self-intersections, and each choice of integral local system $\rho$ on $N$, there are a pair of sets $\Def(N, \rho) $ and $\Obs(N, \rho)$, and a map
%   \begin{equation}
%     \Def(N, \rho) \to \Obs(N, \rho).
%   \end{equation}
%   The space $\Obs(N, \rho) $ is equipped with a canonical element called $0$.
% %  the moduli space of holomorphic discs associated to a compatible almost complex structure $J$, which is conical at infinity, determines a function from $H^1(N)$ to $H^2(N)$. The zero-locus of this function is called the set of bounding cochains.
% \item The triples $(N, \rho, b)$, where $b \in \Def(N, \rho)$ is an element mapping to $0$ in $\Obs(N, \rho)$ are the objects of an  $A_\infty$ category (called the wrapped Fukaya category). We call such a triple a \emph{brane supported on $N$}.
% \item  If $N'$ is Hamiltonian isotopic to $N$, then for each brane supported on $N$, there is a quasi-isomorphic brane supported on $N'$.
% \item The wrapped Fukaya category is generated by the cotangent fibre
% \item The self-Floer cohomology of the cotangent fibre is supported in non-positive degree, and is isomorphic in degree $0$ to the group ring of $\pi_1(Q)$.
%   \item For each representation $\rho$ of $\pi_1(Q)$ which is projectively unitary, there is a brane $\cQ_{\rho}$, supported on a the graph of a closed $1$-form on $Q$, with the property that the Floer cohomology $\cQ_{\rho}$ with a cotangent fibre is isomorphic to $\rho$.
% \end{enumerate}



% These fact assemble to prove the following result:



% \begin{prop}
%   Let $N \subset T^* Q$ be a Maslov-$0$ Lagrangian and $b$ a bounding cochain on $N$, such that the Floer homology $HF^*((N,b), T^*_qQ)$ of $N$ with a cotangent fibre is supported in degree $0$. If the associated representation of $\pi_1(Q)$ admits a non-trivial morphism from the tensor product of an abelian and a unitary representation, then there is an object of the Fukaya category  Floer homology of $(N, \bfb)$ 
% \end{prop}


%  The above definition of isomorphism is essential to formulating the notion of a \emph{generator} of the Fukaya category. In order to formulate it, we recall that the endomorphisms of any fixed object in a linear category form an algebra, and that each other object gives rise to a module over this algebra. More generally, the endomorphisms in an $A_\infty$ category such as the Fukaya category form an $A_\infty$ algebra, and every other object defines an $A_\infty$ module.

%  \begin{dfn} \label{def:generates}
%     An object $(K, \bfb)$ \emph{generates} the Fukaya category if the module $ CF^*( (K, \bfb),(L', \bfb')) $ over $CF^*( (K, \bfb),(K, \bfb)) $ assigned to each object  $(L', \bfb')$ detects isomorphisms, i.e. two objects of the Fukaya category are isomorphic if and only if the associated modules over $CF^*( (K, \bfb),(K, \bfb)) $ are isomorphic in the derived category of modules over $ CF^*( (K, \bfb),(L', \bfb'))$. 
%   \end{dfn}
% The above definition essentially says that, if we can identify a generator of the Fukaya category, and understand the category of $A_\infty$ modules over the endomorphism algebra of this generator, then we can extract information about arbitrary objects of this category.

% In order to apply these ideas to our context, we shall use the 



% given a Riemannian metric on $Q$, we obtain a function
% \begin{equation}
%   r \colon T^* Q \to [0,\infty)
% \end{equation}
% which is given in the standard coordinates by
% \begin{equation}
%   r(p_1, 
% \end{equation}

% Recalling that the symplectic form on the cotangent bundle is the differential of a $1$-form $\lambda$ given in local coordinates by $\sum p_i dq_i$, the contact-type condition is the statement that





  
%   \

% The key 


% \subsection{Curved deformation from holomorphic discs}
% \label{sec:curv-deform-from}



% The key idea of Lagrangian Floer theory is that holomorphic discs determine a deformation of $\Loc(L)$. In order to define it precisely, given a sequence $\{(L_i, E_i)\}_{i=1}^k$ of unitary local systems on Lagrangian branes $L_i$, we let the \emph{classical part} of the graded module $\bigotimes_{i=0}^{k-1}  CF^*(E_i,E_{i+1}) $ consist of the subcomplex which is trivial unless either $L_0 = \cdots = L_{k-1}$ or $L_1 = \cdots = L_k$, in which case it is given by 
% \begin{align}
% &  \Omega^*(L_0 ; \Hom_{\Lambda}(E_0,E_1)) \otimes \cdots \otimes \Omega^*(L_1 ; \Hom_{\Lambda}(E_{k-2},E_{k-1})) \otimes CF^*(E_{k-1}, E_k) \\
% &    CF^*(E_0,E_1) \otimes  \Omega^*(L_1 ; \Hom_{\Lambda}(E_1,E_2)) \otimes \cdots \otimes \Omega^*(L_1 ; \Hom_{\Lambda}(E_{k-1},E_k)),
% \end{align}
% and we take the direct sum if both conditions hold.


% \begin{dfn}
%   A \emph{curved filtered} $A_\infty$ deformation of the classical limit of the Fukaya category consists of a $\Lambda$-linear map
%   \begin{equation}
%     \fm_k :\bigotimes_{i=0}^{k-1}  CF^*(E_i,E_{i+1}) \to   CF^*(E_0,E_{k})
%   \end{equation}
%   for each non-negative integer $0 \leq k$ and sequence $\{(L_i, E_i)\}_{i=1}^k$ of unitary local systems on Lagrangian branes $L_i$ such that:
%   \begin{enumerate}
%   \item The valuation of $\fm_k$ is non-negative, and the reduction modulo the maximal ideal agrees on the classical part with the $A_\infty$ structure induced by the differential graded structure on $\Loc(L)$.
%     \item The operations $\fm_k$ satisfy the $A_\infty$-equation.
%   \end{enumerate}
% \end{dfn}
% To clarify the terminology, the curved structure amounts to the statement that $\fm_0$ may not vanish. The filtered condition is the statement that, if we expand $\fm_k$ as a sum of $\C$-linear maps weighted by $T^\lambda$, there are no contributions for negative $\lambda$, and those for $\lambda = 0$ vanish unless $k=1$ or $k=2$ on the classical part, in which case we obtain the de Rham differential and the wedge product of forms.



% Fixing a Lagrangian $L$ brane as in Definition \ref{def:Lagrangian-brane}, we now state a minor variant of the construction of the Fukaya category:
% \begin{thm} \label{thm:curved_deformation_local_systems}
% A choice of virtual fundamental chains on the moduli spaces of $J$-holomorphic discs with boundary on $L$ determines a curved filtered $A_\infty$-deformation of the classical limit of the Fukaya category. \qed
% \end{thm}
% Replacing de Rham cochains by singular cochains, the construction of this $A_\infty$ deformation generalises the construction of Fukaya, Oh, Ohta, and Ono's book \cite{FOOO}, and is given for immersed Lagrangians in \cite{AJ}. 

%  \begin{rem}
% The construction using differential forms will appear in \cite{AFOOO}. The main difference is that they only consider embedded Lagrangians, and that extract from the virtual fundamental chains on the moduli spaces of holomorphic discs a curved $A_\infty$-deformation of the de Rham algebra $\Omega^*(L; \Lambda)$, without unitary local systems. However, the addition data required to include unitary local systems is discrete, as it consists of keeping track of monodromies of local systems along boundary of holomorphic disc (see, e.g. \cite{Ab4}). Thus, no additional virtual fundamental chain need to be defined.

%  The construction in \cite{AFOOO} is phrased in terms of closed Lagrangians in compact symplectic manifolds. The only difference in the case of manifolds with boundary is that we need to ensure that stable holomorphic discs cannot escape to the boundary. The condition that $J$ be of contact type ensures this (see, e.g. \cite{AbSd}). Finally, we note that the constructions in \cite{FOOO,AJ,AFOOO} produce a \emph{gapped} $A_\infty$ structure, which is a stronger notion than a filtered structure. We omit discussing it since we shall not use it.  
% \end{rem}


% We end this section with a brief discussion of what we mean by a virtual fundamental chain: I DON'T KNOW IF THIS IS WORTH IT!

% \subsection{Bounding cochains and the Fukaya category}
% \label{sec:deform-class-unobstr}

% \begin{dfn}
%   The \emph{deformation and obstruction spaces} of a local system $E$ on a Lagrangian brane are the subspaces
%   \begin{align}
%     \Def(E) & \subset    CF^1(E,E))  \\
%     \Obs(E) & \subset    CF^*(E,E)
%   \end{align}
% consisting of elements of strictly positive valuation.
% \end{dfn}

% The $A_\infty$ structure determines a map
% \begin{align}
%   \Def(L) & \to \Obs(L)  \\
%   b & \mapsto \sum_{k=0}^{\infty} \fm_k(b, \ldots, b).
% \end{align}
% \begin{dfn}
%   A \emph{bounding cochain} is a deformation class $b \in \Def(E)$ whose image in the obstruction space vanishes.
% \end{dfn}

% We now define the Fukaya category $\cF(M)$ to be the $A_\infty$ category with objects such bounding cochains, with morphism spaces
% \begin{equation}
%   CF^*(b_0, b_1) \equiv CF^*(E_0,E_1),
% \end{equation}
% and $A_\infty$ structure given by
% \begin{equation}
%     \fm_k (\phi_0, \ldots, \phi_k) =  \sum_{0 \leq r_i} \mu_{k + \sum r_i}(b_0^{\otimes r_0}, \phi_1, \ldots, b_{k-1}^{\otimes r_{k-1}}, \phi_k, b_k^{\otimes r_k}).
% \end{equation}

% We now list a key property of the Fukaya category:
% \begin{prop} \label{prop:quasi-units}
% Let $\phi$ be a positive Hamiltonian diffeomorphism of $X$ and $b$ a bounding cochain on a local system $E$ supported on a Lagrangian brane $L$. If $b$ a bounding cochain on $E$, there exists a bounding cochain $\phi_* b$ on the local system induced by $E$ on $\phi L$, and an element
% \begin{equation}
%   \kappa_{\phi} \in HF^0(b, \phi_* b)  
% \end{equation}
% depending only on the induced path of Lagrangians, such that, for any bounding cochain $b'$ supported on a compact Lagrangian brane, the induced maps:
% \begin{align}
%      HF^*(b',b) & \to   HF^*(b', \phi_* b) \\
%  HF^*(\phi_* b,b) & \to  HF^*(b,b')
% \end{align}
% are isomorphisms.
% \end{prop}
% \begin{proof}s
% The construction of $\phi_* b$ and of the element $\kappa_{\phi}$ given in \cite{AJ} for closed Lagrangians applies to Lagrangians with Legendrian boundary as long as we assume positivity of the Hamiltonian isotopy, which ensure the the moduli spaces are compact. To prove $\kappa_\phi$ induces an isomorphism with respect to Floer cohomology of compactly supported branes, we use the canonical isomorphism $HF^*(b', \phi_* b) = HF^*(\phi_*^{-1} b', b) $, and the fact that any Hamiltonian isotopy of compact Lagrangians is isotopic to one which is supported away from the boundary.
% \end{proof}


% \begin{dfn}
% A Lagrangian $L$ with Legendrian boundary is {\it unobstructed} if it supports a brane, a local system, and a bounding cochain.
% \end{dfn}





\begin{thebibliography}{999}

%\bibitem{AbSm} M Abouzaid and I Smith
%`Exact Lagrangians in plumbings'
%Geom. Funct. Anal. 22 (2012) 785--831

\bibitem{AJ} M Akaho and D Joyce
`Immersed Lagrangian Floer theory' J.\ Differential Geom.\ 86 (2010) 381--500


\bibitem{Fuk1} K Fukaya
`Cyclic symmetry and adic convergence in Lagrangian Floer theory'
Kyoto J.\ Math.\ 50 (2010) 521--590. 

\bibitem{Fuk2} K Fukaya
`Unobstructed immersed Lagrangian correspondence and filtered $A_\iy$ functor'
arXiv:1706.02131

\bibitem{FOOO} K Fukaya, Y-G Oh, H Ohta and K Ono
`Lagrangian intersection Floer theory---anomaly and obstruction'
Parts I and II AMS/IP Studies in Advanced Mathematics 46.1 and 46.2 (2009)
%and Chapter 10
%at \url{https://www.math.kyoto}-\url{u.ac.jp/~fukaya/fukaya.html}

\bibitem{FOOOt} K Fukaya, Y-G Oh, H Ohta and K Ono
`Lagrangian Floer theory on compact toric manifolds. I' Duke Math.\ J.\ 151 (2010) 23--174

%`Lagrangian Floer theory and mirror symmetry on compact toric manifolds. I'
%Asterisque, vol 376, 2016, Societe Mathematique de France

\bibitem{FOOO-Z} K Fukaya, Y-G Oh, H Ohta and K Ono
'Lagrangian Floer theory over integers: spherically positive symplectic manifolds'
Pure Appl.\ Math.\ Q.\
9 (2013) 189--289

\bibitem{HL} R Harvey and H B Lawson Jr
`Calibrated geometries' Acta Math.\ 148 (1982) 47--157

\bibitem{IJO} Y Imagi, D Joyce and J Oliveira dos Santos
`Uniqueness results for special Lagrangians and Lagrangian mean curvature flow in $\C^m$'
Duke Math.\ J 165 (2016) 847--933

\bibitem{Jb} D Joyce
`Riemannian holonomy groups and calibrated geometry'
Oxford Graduate Texts in Mathematics 12 (2007) 

\bibitem{Jc} D Joyce
`Conjectures on Bridgeland stability for Fukaya categories of Calabi--Yau manifolds, special Lagrangians, and Lagrangian mean curvature flow'
EMS Surv.\ Math.\ Sci.\ 2 (2015) 1--62


\bibitem{KS} K Kodaira and D C Spencer
`On deformations of complex analytic structures, III. Stability theorems for complex structures'
Ann.\ of Math.\ 71 (1960) 43--76


\bibitem{Lj} S \L{}ojasiewicz
`Ensembles semi-analytiques'
Inst.\ Hautes \'Etudes Sci.\ 1965



\bibitem{LP}
J D Lotay and T Pacini
`From minimal Lagrangian to J-minimal submanifolds: persistence and uniqueness'
Boll.\ Unione Mat.\ Ital.\ 12 (2019) 63--82
\bibitem{LP2}
J D Lotay and T Pacini
`From Lagrangian to totally real geometry: coupled flows and calibrations' Comm.\ Anal.\ Geom.\ (2020) 607--675

\bibitem{ML} R C McLean
`Deformations of calibrated submanifolds' Comm.\ Anal.\ Geom.\ 6 (1998) 705--747


\bibitem{Sd1} P Seidel
`Graded Lagrangian submanifolds'
Bull.\ Soc.\ Math.\ Fr.\ 128 (2000) 103--149

\bibitem{TY} R P Thomas and S-T Yau
`Special Lagrangians, stable bundles and mean curvature flow'
Comm.\ Anal.\ Geom.\ 10 (2002) 1075--1113

\end{thebibliography}
\end{document}
%%% Local Variables:
%%% mode: latex
%%% TeX-master: t
%%% End:


\bibitem{Fk1} K Fukaya
`Morse homotopy, $A^\iy$-category, and Floer homologies' Proc. GARC Workshop (1993) 1--102

\bibitem{Fk2} K Fukaya
`Floer homology and mirror symmetry. II' Adv. Stud. Pure Math. 34 (2002) 31--127

\bibitem[FOOO09]{FOOO} K Fukaya, Y-G Oh, H Ohta and K Ono
`Lagrangian Intersection Floer Theory---Anomaly and Obstruction'
Parts I \& II AMS/IP 46.1 \& 46.2 (2009)
and Chapter 10
at \url{https://www.math.kyoto}-\url{u.ac.jp/~fukaya/fukaya.html}

\bibitem[FOOO10]{FOOO2} K Fukaya, Y-G Oh, H Ohta and K Ono
`Lagrangian Floer theory on compact toric manifolds. I' Duke Math. J. 151 (2010) 23--174

\bibitem[FOOO13]{FOOO-Z} K Fukaya, Y-G Oh, H Ohta and K Ono
'Lagrangian Floer theory over integers: spherically positive
   symplectic manifolds'
Pure Appl. Math. Q.
9 (2013) 189--289

\bibitem{FSS1} K Fukaya, P Seidel and I Smith
`Exact Lagrangian submanifolds in simply-connected cotangent bundles'
Invent. Math. 172 (2008) 1--27

% \bibitem{FSS2} K Fukaya, P Seidel and I Smith
% `The symplectic geometry of cotangent bundles from a categorical viewpoint'
% Lect. Notes Phys. 757 (2009) 1--26

%\bibitem{GS} V Guillemin and M Stenzel
%`Grauert tubes and the homogeneous Monge--Amp\`ere Equation'
%J. Diff. Geom. 34 (1991) 561--570

%\bibitem{GT} D Glibarg and N S Trudinger
%`Elliptic Partial Differential Equations of Second Order'
%Grund. Math. Wiss. 224 (1977)

%\bibitem{dG} E de Giorgi
%`Frontiere orientate di misura minima' Editr. Tecn. Scient. 1961

\bibitem{Gr} H Grauert
`On Levi's problem and the imbedding of real-analytic manifolds'
Ann. of Math. 68 (1958) 460--472

%\bibitem{Gr} M Gromov
%`Pseudoholomorphic curves in symplectic manifolds' Invent. Math. 82 (1985) 307--347

\bibitem{GPS}

\bibitem{HL} R Harvey and H B Lawson Jr
`Calibrated geometries' Acta Math. 148 (1982) 47--157

%\bibitem{H1} N J Hitchin
%`Monopoles and geodesics' Comm. Math. Phys. 83 (1982) 579--602

\bibitem{Ht} N J Hitchin
`Stable bundles and integrable systems' Duke Math. J 54 (1987) 91--114

%\bibitem{H3} N J Hitchin
%`A note on vanishing theorems' in `Geometry and Analysis on Manifolds'
%Progr. Math. 308 (2015) 373--382

\bibitem{Hn} K Honda
`A note on Morse theory of harmonic 1-forms'
Topol. 38 (1999) 223--233

\bibitem{Hp} H Hopf
`Vektorfelder in $n$-dimensionalen Mannigfaltigkeiten'
Math. Annal. 96 (1926) 225--250

\bibitem{Hr} H Hironaka
`Resolution of singularities of an algebraic variety over a field of characteristic zero'
Ann. of Math. 79 (1964) 109--203 205--326

%\bibitem{I} Y Imagi
%`Surjectivity of a gluing construction in special Lagrangian geometry'
%to appear in Comm. Anal. Geom.

\bibitem{IJO} Y Imagi, D Joyce and J Oliveira dos Santos
`Uniqueness results for special Lagrangians and Lagrangian mean curvature flow in $\C^m$'
Duke Math. J 165 (2016) 847--933

%\bibitem{J0} D Joyce
%`On counting special Lagrangian homology 3-spheres'
%Contemp. Math. 314 (2002) 125--151

\bibitem{J5} D Joyce
`Special Lagrangian submanifolds with isolated conical singularities V. Survey and applications'
J Diff. Geom. 63 (2003) 279--347

%\bibitem{Jc} D Joyce
%`Riemannian Holonomy Groups and Calibrated Geometry' OUP 2006

%\bibitem{JLT} D Joyce, Y-I Lee and M-P Tsui
%`Self-similar solutions and translating solitons for Lagrangian mean curvature flow'
%J. Diff. Geom. 84 (2010) 127--161

%\bibitem{Jc} D Joyce
%`Conjectures on Bridgeland stability for Fukaya categories of Calabi--Yau manifolds, special Lagrangians, and Lagrangian mean curvature flow'
%EMS Surv. Math. Sci. 2 (2015) 1--62

\bibitem{Kl} E K\"ahler
`Einf\"uhrung in die Theorie der Systeme von Differentialgleichungen'
B G Teubner 1934

\bibitem{Kld} D Kaledin and M Verbitsky
`Hyperkahler Manifolds' Int. Press 1999

%\bibitem{Kp} N Kapouleas
%`Doubling and desingularization constructions for minimal surfaces'
%Adv. Lect. Math. 20 (2011) 281--325

%\bibitem{Kp2} N Kapouleas
%`Minimal surfaces in the round three-sphere by doubling the equatorial two-sphere, I'
%arXiv:1409.0226

\bibitem{KY} N Kapouleas and S-D Yang
`Minimal surfaces in the three-sphere by doubling the Clifford torus'
Amer. J Math. 132 (2010) 257--295

%\bibitem{Kb1} S Kobayashi
%`First Chern class and holomorphic tensor fields'
%Nag. Math. J 77 (1980) 5--11

%\bibitem{Kb2} S Kobayashi
%`Curvature and stability of vector bundles'
%Proc. Jap. Acad. 58 (1982) 158--162

\bibitem{Kn} M Kontsevich
`Homological algebra of mirror symmetry' Proc. ICM 1994

\bibitem{Krause}
  H Krause, 'Localization theory for triangulated categories' 
Triangulated categories
  London Math. Soc. Lecture Note Ser. 375 (2010)
161--235
%\bibitem{KM} P B Kronheimer and T Mrowka
%`Monopoles and three-manifolds' New Mathematical Monographs 10 (2007)

\bibitem{Ll} G Lawlor
`The angle criterion' Invent. Math. 95 (1989) 437--446

%\bibitem{Lee} Y-G Lee
%`Embedded special Lagrangian submanifolds in Calabi--Yau manifolds'
%Comm. Anal. Geom. 11 (2003) 391--423

%\bibitem{LS1} C de Lellis and E Spadaro
%`Q-valued functions revisited' Mem. AMS 211 (2011)

\bibitem{LS} C De Lellis and E Spadaro
`Regularity of area minimizing currents I: gradient $L^p$ estimates'
Geom. Funct. Anal. 24 (2014) 1831--1884

%`Regularity of area minimizing currents II \& III'
%gradient $L^p$ estimates,
%center manifold and blow-up'
%Geom. Funct. Anal. 24 (2014) 1831--1884
%Ann. Math. 183 (2016) 499--575 \& 577--617


%\bibitem{LS3} C de Lellis and E Spadaro
%`Multiple valued functions and integral currents' Ann. Sc. Norm. Sup. Pisa 14 (2015) 1239--1269
 

%\bibitem{LS} Lempert and Sz\"oke

\bibitem{Lie} S Lie
`Theorie des Transformations-Gruppen'
Arch. Math. Nat. 1 (1876) 152--193

\bibitem{Lj} S \L{}ojasiewicz
`Ensembles semi-analytiques'
IHES 1965

\bibitem{LP}
%\bibitem{MW} M Marcolli and B-L Wang
%`Equivariant Seiberg-Witten Floer homology' Comm. Anal. Geom. 9 (2001) 451--639

%\bibitem{MS} D Mcduff and D Salamon
%`$J$-holomorphic curves and symplectic topology' 2nd edition AMS 2012

\bibitem{ML} R C McLean
`Deformations of calibrated submanifolds' Comm.\ Anal.\ Geom.\ 6 (1998) 705--747

\bibitem{Miln}
J W Milnor h-cobordism

\bibitem{Mr} C B Morrey Jr
`Multiples Integrals in the Calculus of Variations'
Grund. Math. Wiss. 130 (1966)

\bibitem{Ms} J Moser
`On the volume elements on a manifold'
Trans. Amer. Math. Soc. 120 (1965) 286--294

%\bibitem{Nsh1} J Nash
%`Real algebraic manifolds'
%Ann. Math. 56 (1952) 405--421

%\bibitem{Nsh2} J Nash
%`The imbedding problem for Riemannian manifolds'
%Ann. Math. 63 (1956) 20--63

%\bibitem{Nev1} A Neves
%`Singularities of Lagrangian mean curvature flow: zero-Maslov class case' Invent. Math. 168  (2007) 449--484
 
%\bibitem{Nev2} A Neves
%`Recent progress on singularities of Lagrangian mean curvature flow'
%Adv. Lect. Math. 20 (2011) 413--438

\bibitem{Nv} S P Novikov
`Multivalued functions and functionals. An analogue of the Morse theory'
Sov. Math. Dukl 24 (1981)

\bibitem{Oh}
  Y-G Oh,
   'Seidel's long exact sequence on Calabi--Yau manifolds'
Kyoto J. Math. 51 (2011) 687--765
%\bibitem{OS} P Ozsv\'ath and Z Szab\'o
%`Holomorphic triangle invariants and the topology of symplectic four-manifolds'
%Duke Math. J 121 (2004) 1--34

%\bibitem{Pzh} A N Pazhitnov
%`On the Novikov complex for rational Morse forms'
%Ann. Fac. Sci. Toul. 4 (1995) 297--338

%\bibitem{PR} J T Pitts and J H Rubinstein
%`Equivariant minimax and minimal surfaces in geometric three-manifolds'
%Bull. AMS 19 (1988) 303--309

%\bibitem{Pl} L Polterovich
%`Surgery of Lagrange submanifolds'
%Geom. Funct. Anal. 1 (1991) 198--210

%\bibitem{SU} Sacks


\bibitem{Sd1} P Seidel
`Graded Lagrangian submanifolds'
Bull. Soc. Math. Fr. 128 (2000) 103--149

\bibitem{Sd2} P Seidel
`A long exact sequence for symplectic Floer cohomology'
Topol. 42 (2003) 1003--1063

\bibitem{Sd3} P Seidel
`Fukaya Categories and Picard--Lefschetz Theory' EMS 2008

\bibitem{Seidel} Seidel
%\bibitem{Sd4} P Seidel
%`Homological Mirror Symmetry for the Quartic Surface'
%Mem. AMS 236 (2015)

%\bibitem{Sh} N Sheridan
%`Homological mirror symmetry for Calabi--Yau hypersurfaces in projective space' Invent. Math. 199 (2015) 1--186

%\bibitem{Sm1} L Simon
%`Asymptotics for a class of nonlinear evolution equations, with applications to geometric problems'
%Ann. Math. 118 (1983) 525--571

\bibitem{Sm}
L Simon `Lectures on Geometric Measure Theory' ANU 1983

\bibitem{St} M B Stenzel
`Ricci-flat metrics on the complexification of a compact rank one symmetric
space' Manus. Math. 80 (1993) 151--163

%\bibitem{SW1} N Seiberg and E Witten
%`Electric-magnetic duality, monopole condensation, and confinement in N=2 supersymmetric Yang--Mills theory' Nucl. Phys. B 426 (1994) 19--52

%\bibitem{SW2} N Seiberg and E Witten
%`Monopoles, duality and chiral symmetry breaking in N=2 supersymmetric QCD' Nucl. Phys. B 431 (1994) 484--550

%\bibitem{Tb} C H Taubes
%`Self-dual Yang--Mills connections on non-self-dual 4-\0{}manifolds'
%J Diff. Geom. 17 (1982) 139--170

%\bibitem{Tb2} C H Taubes
%`The Seiberg-Witten equations and the Weinstein conjecture' Geom. Topol. 11 (2007)  2117--2202

%\bibitem{Tm} R P Thomas
%`Moment maps, monodromy and mirror manifolds'
%in `Symplectic geometry and mirror symmetry' World Sci. (2001) 467--498

%\bibitem{Tm2} R P Thomas
%`A holomorphic Casson invariant for Calabi--Yau 3-folds, and bundles on K3
%fibrations'
%J. Diff. Geom. 54 (2000) 367--438

\bibitem{TY} R P Thomas and S-T Yau
`Special Lagrangians, stable bundles and mean curvature flow'
Comm. Anal. Geom. 10 (2002) 1075--1113

\bibitem{Tits} J Tits
`Free subgroups in linear groups'
J Alg. 20 (1972) 250--270

%\bibitem{Tg} A Tognoli
%`Su una congettura di Nash' Ann. Sc. Norm. Sup. Pisa 27 (1973) 167--185

\bibitem{TW} C-J Tsai and M-T Wang.
`A strong stability condition on
minimal submanifolds and its implications'
J. reine angew. Math, ahead of print, published online at
\url{https://doi.org/10.1515/crelle-2018-0038}

%`Mean curvature flows in manifolds of special holonomy'
%to appear in J. Diff. Geom.

%\bibitem{Ul}
%\bibitem{UY} K Uhlenbeck and S-T Yau
%`On the existence of Hermitian--Yang--Mills connections in stable vector bundles'
%Comm. Pure Appl. Math. 39 (1986) 257--293 

\bibitem{Wn} A Weinstein
`Symplectic manifolds and their Lagrangian submanifolds'
Adv. Math. 6 (1971) 329--346

%\bibitem{Wtt} E Witten
%`Two-dimensional gravity and intersection theory on moduli space'
%Surv. Diff. Geom. (1991) 243--310

%\bibitem{Wt} Whitney

%\bibitem{Witt2} E Witten
%`Monopoles and four-manifolds' Math. Res. Lett. 1 (1994) 769--796 

%\bibitem{YM} Yang--Mills

%\bibitem{Y} S-T Yau
%`On the Ricci curvature of a compact K\"ahler manifold and the complex Monge--Amp\`ere equation, I'
%Comm. Pure Appl. Math. 31 (1978) 339--411
\end{thebibliography}
\end{document}
%%% Local Variables:
%%% mode: latex
%%% TeX-master: t
%%% End:
